\documentclass[11pt,a4paper]{article}

% ── Packages ──
\usepackage[utf8]{inputenc}
\usepackage[T1]{fontenc}
\usepackage{amsmath,amssymb,amsthm}
\usepackage{mathtools}
\usepackage{algorithm}
\usepackage{algpseudocode}
\usepackage{booktabs}
\usepackage{graphicx}
\usepackage[margin=2.5cm]{geometry}
\usepackage{hyperref}
\usepackage{cleveref}
\usepackage{natbib}
\usepackage{enumitem}
\usepackage{subcaption}
\usepackage{xcolor}
\usepackage{float}
\usepackage{multirow}

\providecommand{\doi}[1]{\href{https://doi.org/#1}{\texttt{doi:#1}}}

\newtheorem{definition}{Definition}
\newtheorem{proposition}{Proposition}

\title{Optimal Route Planning for Orienteering Races on Real Terrain:\\
An Integrated Cost Surface and Branch-and-Cut Approach}

\author{
  Pablo Borrego Ramos\thanks{Independent Researcher.}
}

\date{\today}

\begin{document}
\maketitle

\begin{abstract}
We present an end-to-end computational framework for provably optimal route planning in orienteering races on real terrain.
A GIS pipeline combines IOF map symbols, DEM-derived morphometry, and the Minetti metabolic model into a fully asymmetric cost matrix, where directional slope costs and cumulative fatigue make travel cost depend on both direction and elapsed effort.
The terrain realities captured by this pipeline---asymmetry, fatigue, and heterogeneous surface costs---lead naturally to a new problem variant, the Asymmetric Orienteering Problem with Fatigue (AOPF), for which we develop tailored exact solution methods.
We propose two LP formulations: an MTZ-based model with McCormick linearisation and a single-commodity flow formulation where the fatigue budget is natively linear.
We strengthen the Branch-and-Cut solver with routing infeasibility cuts, a novel family of pure node-variable inequalities based on the assignment relaxation, alongside tightened flow-arc coupling bounds, fatigue-aware lifted covers, and directed adaptations of cycle cover and path inequalities.
We validate the framework on two Spanish terrains across seven control-point distributions each, and release a parameterised benchmark suite of 21 instances for the asymmetric OP with fatigue.
The flow-based B\&C certifies global optimality on 9 of 14 real-terrain instances within a 15-minute time limit, with tighter LP gaps than the MTZ formulation (mean gap 5.4\% vs.\ 16.0\% on benchmark instances).
\end{abstract}

\noindent\textbf{Keywords:} orienteering problem, asymmetric routing, cumulative fatigue, branch-and-cut, valid inequalities, flow formulation, terrain cost surface, Minetti metabolic model

% ════════════════════════════════════════════════════════════════
\section{Introduction}
\label{sec:intro}
% ════════════════════════════════════════════════════════════════

The Orienteering Problem (OP) asks a competitor to select a subset of control points and determine a route starting and ending at a depot, maximising the total score collected while respecting a travel budget.
Formally introduced by \citet{tsiligirides1984heuristic}, the OP is NP-hard and has been studied extensively in operations research.
However, most formulations assume symmetric, Euclidean distances, which poorly approximate real-world orienteering where terrain, vegetation, slope direction, and fatigue create highly asymmetric travel costs.

In competitive orienteering and rogaine events, route choice is the decisive skill: the runner must balance point value against travel difficulty across heterogeneous terrain.
Course setters estimate ideal times using subjective experience, but the underlying cost function is a nonlinear function of terrain type, slope, and accumulated fatigue.

This paper bridges two research streams: (i)~GIS-based cost surface modelling for orienteering maps, following \citet{ica2021lcp}, and (ii)~exact optimisation of the OP via integer programming.
The terrain realities captured by the cost pipeline---directional slope costs, surface heterogeneity, and cumulative fatigue---produce a problem structure not addressed by existing OP formulations, motivating both a new problem variant and tailored solution methods.
Our contributions are:
\begin{enumerate}[nosep]
  \item A terrain cost model combining IOF map symbols (ACR), DEM-derived morphometry (HCR), directional metabolic cost (Minetti), and cumulative fatigue into a single asymmetric cost matrix.
  \item An IQR-based scaling method for merging ACR and HCR that adapts to terrain characteristics without manual calibration.
  \item The Asymmetric Orienteering Problem with Fatigue (AOPF), a problem variant arising naturally from the terrain pipeline, formalising the OP with directed travel costs and a cumulative fatigue multiplier.
  To the best of our knowledge, this is the first exact formulation for the OP with position-dependent fatigue costs.
  \item Two LP formulations for the AOPF: (a)~an MTZ-based model with McCormick linearisation of the bilinear fatigue term, and (b)~a single-commodity flow formulation where the fatigue budget is natively linear, eliminating the McCormick relaxation gap entirely.
  The MTZ formulation serves as a natural baseline with fewer variables ($O(n)$ time variables vs $O(n^2)$ flow variables); comparing the two isolates the McCormick envelope as the primary source of LP relaxation looseness.
  \item Valid inequalities tailored to the AOPF: (a)~tightened flow-arc coupling bounds exploiting minimum arrival times, (b)~fatigue-aware lifted cover inequalities with position-dependent cost inflation, and (c)~routing infeasibility cuts---a novel family of pure node-variable inequalities that use the assignment relaxation to identify subsets of nodes that cannot all be visited within the fatigue-adjusted budget.
  We also adapt the cycle cover and path inequalities of \citet{fischetti1998solving} to the directed asymmetric setting with fatigue.
  \item A parameterised benchmark suite of 21 synthetic instances for the asymmetric OP with fatigue, varying node count, asymmetry, fatigue rate, and budget tightness---the first benchmark for this problem variant.
  \item Computational experiments on two real terrains comparing the MTZ and flow formulations, and an ablation study on the benchmark suite quantifying the marginal contribution of each cut family.
\end{enumerate}

The remainder of this paper is organised as follows.
\Cref{sec:related} reviews related work.
\Cref{sec:cost} describes the cost surface model.
\Cref{sec:formulation} presents the mathematical formulation.
\Cref{sec:methods} details the solution methods.
\Cref{sec:pipeline} outlines the computational pipeline.
\Cref{sec:experiments} describes the experimental setup, and \cref{sec:results} presents the results, including a synthetic benchmark suite for the asymmetric OP with fatigue.
\Cref{sec:discussion} discusses findings and limitations, and \cref{sec:conclusion} concludes.

% ════════════════════════════════════════════════════════════════
\section{Related Work}
\label{sec:related}
% ════════════════════════════════════════════════════════════════

\paragraph{The Orienteering Problem.}
The OP was introduced by \citet{tsiligirides1984heuristic} and has since generated a substantial body of literature.
\citet{vansteenwegen2011orienteering} provide a comprehensive survey covering exact methods, heuristics, and variants including the Team OP, OP with Time Windows, and stochastic OP.
Exact approaches typically employ branch-and-bound or branch-and-cut on integer programming formulations \citep{fischetti1998solving,gendreau1998generalized}.
Most formulations assume symmetric distances derived from Euclidean or road-network metrics.
\citet{gunawan2016orienteering} extend the survey to recent variants and applications, noting that terrain-aware formulations remain scarce.
Most recently, \citet{kobeaga2021revisited} proposed a revisited B\&C algorithm for large-scale symmetric OP instances (up to 7{,}397 nodes), incorporating shrinking techniques for subtour elimination, blossom separation heuristics for comb inequalities, efficient variable pricing, and a structured separation loop.
Their component evaluation methodology---systematically enabling and disabling algorithmic components to measure marginal contributions---informs the ablation study in the present work.
Our formulation differs fundamentally in addressing asymmetric costs and cumulative fatigue, which preclude the use of undirected comb inequalities and require the flow-based linearisation developed here.

\paragraph{Metaheuristics for the OP.}
Given the NP-hardness of the OP, metaheuristics are widely used for practical instances.
Simulated Annealing \citep{kirkpatrick1983sa}, Genetic Algorithms \citep{tasgetiren2002ga_op}, and Ant Colony Optimisation \citep{ke2008aco_op} have all been applied.
\citet{vansteenwegen2009iterated} propose an iterated local search that is competitive with exact methods on benchmark instances.
Our SA implementation incorporates Or-opt moves and stagnation detection, with parameters tuned via grid search over 1{,}000 random instances.

\paragraph{Terrain cost modelling.}
Least-cost path (LCP) analysis on terrain has a long history in GIS.
\citet{tobler1993} proposed the hiking function relating slope to walking speed, which remains widely used despite known limitations on steep descents.
\citet{minetti2002} measured the metabolic cost of locomotion across a wide range of slopes, producing a fifth-degree polynomial that better captures the energetics of both ascent and descent.
\citet{rees2004} compared several slope-speed models and found significant differences in predicted travel times.

\paragraph{Orienteering map cost surfaces.}
\citet{ica2021lcp} introduced a GIS-based approach to route planning on orienteering maps, assigning travel speeds to IOF symbol categories and computing least-cost paths.
They proposed the Hypsometry Cost Raster (HCR) combining terrain ruggedness, plan curvature, and slope magnitude to capture morphometric difficulty not represented by map symbols.
Our work extends their approach by: (i)~replacing Tobler's hiking function with the Minetti metabolic model for directional slope costs, (ii)~proposing IQR-based scaling instead of histogram mode matching for ACR-HCR fusion, and (iii)~coupling the cost surface with an exact OP solver rather than point-to-point LCP analysis.

\paragraph{Asymmetric and state-dependent routing.}
The Asymmetric Travelling Salesman Problem (ATSP) has a rich polyhedral theory, including directed subtour elimination constraints and facet-defining inequalities such as the $D_k^+$ and odd-CAT families \citep{fischetti1991atsp,fischetti1997atsp}.
Our directed SECs and connectivity cuts follow the standard ATSP formulation; the budget-based cuts (B1--B4) are specific to the selective (orienteering) setting where node visits are optional.
The fatigue model in the AOPF makes arc costs state-dependent: the effective cost of traversing an arc depends on the cumulative cost elapsed along the route.
State-dependent travel costs arise in several routing contexts, including load-dependent vehicle routing \citep{bektas2011pollution}, time-dependent VRP \citep{malandraki1992timedep}, and energy-minimising routing \citep{kara2007energy}.
In these problems, the bilinear interaction between a state variable and arc selection creates linearisation challenges analogous to our fatigue term; our flow-based formulation addresses this without McCormick envelopes.
The OP is a member of the broader family of TSP with Profits \citep{feillet2005tsp_profits}, which includes the prize-collecting TSP and the profitable tour problem.
To the best of our knowledge, no prior work in this family addresses position-dependent cumulative fatigue costs with exact methods.

% ════════════════════════════════════════════════════════════════
\section{Cost Surface Model}
\label{sec:cost}
% ════════════════════════════════════════════════════════════════

Let the terrain be discretised into a regular grid of cells with resolution $\delta$ (metres per cell), where each cell is indexed by its row $r$ and column $c$.
Each cell $(r,c)$ is assigned a base traversal cost that integrates map-derived and DEM-derived information.

% ── 2.1 ACR ──
\subsection{Areal Cost Raster (ACR)}
\label{sec:acr}

The ACR assigns a dimensionless weight $w_{rc} \in (0, 10]$ to each cell based on the IOF map symbol classification.
Following \citet{ica2021lcp}, each symbol category maps to a running speed, and the weight is inversely proportional:
\begin{equation}
  w_{rc} = \frac{\delta}{v_{\text{base}} \cdot s_{\text{cat}(r,c)}}
  \label{eq:acr}
\end{equation}
where $v_{\text{base}} = 2.5$\,m/s is the reference flat-terrain speed, $s_{\text{cat}}$ is the speed factor for the symbol category (e.g., $s = 1.0$ for paved road, $s = 0.133$ for fight vegetation), and $\delta$ is the cell size.
Impassable features (water, cliffs, buildings) receive $w = 10$.

The ACR is rasterised from the orienteering map file (OMAP format) by parsing symbol objects---lines, areas, and points---and rendering them onto the UTM grid with appropriate buffer widths.
B\'{e}zier curves are tessellated to preserve cartographic accuracy.

% ── 2.2 HCR ──
\subsection{Hypsometry Cost Raster (HCR)}
\label{sec:hcr}

The HCR captures terrain morphometry not represented by map symbols, derived from a Digital Terrain Model (DTM).
Following \citet{ica2021lcp}, three morphometric variables are computed and normalised to $[\ell, u] = [0.1, 1.0]$:

\paragraph{Terrain Ruggedness Index (TRI).}
The standard deviation of elevation in a $3 \times 3$ neighbourhood \citep{riley1999tri}:
\begin{equation}
  \text{TRI}_{rc} = \sqrt{\frac{1}{9}\sum_{(i,j) \in \mathcal{N}(r,c)} \bigl(z_{ij} - \bar{z}_{rc}\bigr)^2}
  \label{eq:tri}
\end{equation}
where $z_{ij}$ is the elevation at cell $(i,j)$ and $\bar{z}_{rc}$ is the local mean.

\paragraph{Plan Curvature (PC).}
The absolute value of plan curvature detects ridges, spurs, and side valleys \citep{zevenbergen1987}:
\begin{equation}
  \text{PC}_{rc} = \left| \frac{-\bigl(p^2 r - 2pqs + q^2 t\bigr)}{\bigl(p^2 + q^2\bigr)^{3/2}} \right|
  \label{eq:pc}
\end{equation}
where $p = \partial z / \partial x$, $q = \partial z / \partial y$, $r = \partial^2 z / \partial x^2$, $s = \partial^2 z / \partial x \partial y$, $t = \partial^2 z / \partial y^2$, computed via central finite differences at resolution $\delta$.

\paragraph{Slope Magnitude (SLO).}
The slope angle in degrees:
\begin{equation}
  \text{SLO}_{rc} = \arctan\!\sqrt{\left(\frac{\partial z}{\partial x}\right)^2 + \left(\frac{\partial z}{\partial y}\right)^2}
  \label{eq:slo}
\end{equation}

Each variable is normalised via robust percentile clipping:
\begin{equation}
  \hat{V}_{rc} = \ell + (u - \ell) \cdot \frac{\text{clip}(V_{rc}, V^{(1\%)}, V^{(99\%)}) - V^{(1\%)}}{V^{(99\%)} - V^{(1\%)}}
  \label{eq:norm}
\end{equation}

The HCR is then:
\begin{equation}
  \text{HCR}_{rc} = \widehat{\text{TRI}}_{rc}^{\,\alpha} \cdot |\widehat{\text{PC}}_{rc}|^{\,\beta} \cdot \widehat{\text{SLO}}_{rc}^{\,\gamma}
  \label{eq:hcr}
\end{equation}
with $\alpha = 1$, $\beta = 0.5$, $\gamma = 1$ as specified by \citet{ica2021lcp}.

% ── 2.3 IQR Scaling ──
\subsection{IQR-Based Raster Fusion}
\label{sec:iqr}

To combine ACR and HCR, the HCR must be scaled into ACR's numerical space.
\citet{ica2021lcp} proposed histogram mode matching: $c = f_{\text{ACR}}(\mu) / f_{\text{HCR}}(\mu)$, where $f(\mu)$ denotes the density at the mode.
This is sensitive to distribution shape and can produce unstable scaling constants.
For example, when the ACR distribution is multimodal (due to discrete symbol categories) and the HCR distribution is unimodal and right-skewed, the mode densities $f(\mu)$ can differ by an order of magnitude, yielding scaling constants that vary erratically with small changes in bin width.

We evaluated six scaling methods on both study areas: median ratio, mean ratio, IQR ratio, standard deviation ratio, percentile range matching, and full quantile mapping.
The median and mean ratios produced excessively large constants ($c > 7$) due to the different distribution shapes of ACR (discrete symbol weights) and HCR (continuous, right-skewed).
Percentile and quantile mapping over-corrected by forcing HCR to mimic ACR's distribution.
The IQR ratio was selected for its robustness to outliers and its focus on matching variability rather than central tendency, adapting automatically to terrain characteristics ($c = 0.87$ for Torremocha, $c = 2.18$ for La Muela).
A straightforward alternative, normalising both rasters to $[0,1]$, was rejected because it would discard the ACR's physically meaningful scale, where cell weights directly encode running speed ratios derived from IOF symbol specifications.

We propose an Interquartile Range (IQR) ratio:
\begin{equation}
  c = \frac{\text{IQR}(\text{ACR})}{\text{IQR}(\text{HCR})} = \frac{Q_{75}^{\text{ACR}} - Q_{25}^{\text{ACR}}}{Q_{75}^{\text{HCR}} - Q_{25}^{\text{HCR}}}
  \label{eq:iqr}
\end{equation}
where $Q_p$ denotes the $p$-th percentile over valid cells.
This matches the \emph{variability} of HCR to ACR rather than aligning central tendencies, producing a more stable and interpretable scaling.

The scaled HCR is added to ACR with path protection:
\begin{equation}
  \text{Cost}_{rc} = w_{rc} + c \cdot \text{HCR}_{rc} \cdot \pi_{rc}
  \label{eq:combined}
\end{equation}
where the path protection factor $\pi_{rc}$ prevents HCR from distorting well-mapped features:
\begin{equation}
  \pi_{rc} = \begin{cases}
    0.1 & \text{if } w_{rc} < 0.15 \text{ (paths/roads)} \\
    0.3 & \text{if } 0.15 \leq w_{rc} < 0.25 \text{ (good tracks)} \\
    0.0 & \text{if } w_{rc} \geq 5.0 \text{ (impassable)} \\
    1.0 & \text{otherwise}
  \end{cases}
  \label{eq:pathprot}
\end{equation}

The path protection factor attenuates HCR on cells where the cartographer has already provided precise cost information.
Paths and roads (ISOM 501--505) are mapped with high positional accuracy in orienteering maps; DEM-derived roughness at these locations reflects the surrounding terrain rather than the path surface.
The attenuation values (0.1 for paths, 0.3 for tracks) were calibrated to allow minimal HCR influence on linear features while preserving full contribution on open terrain, where map symbols provide only coarse vegetation categories.
Impassable features ($w \geq 5.0$) receive zero HCR contribution to prevent softening of hard barriers.

% ── 2.4 Minetti ──
\subsection{Directional Metabolic Cost (Minetti Model)}
\label{sec:minetti}

The metabolic cost of locomotion on a slope gradient $i$ (rise/run, dimensionless) is modelled by the fifth-degree polynomial of \citet{minetti2002}:
\begin{equation}
  C(i) = 280.5\,i^5 - 58.7\,i^4 - 76.8\,i^3 + 51.9\,i^2 + 19.6\,i + 2.5 \quad [\text{J/kg/m}]
  \label{eq:minetti}
\end{equation}
valid for $i \in [-0.45, +0.45]$.
The cost multiplier relative to flat terrain is:
\begin{equation}
  \phi(i) = \frac{\max\bigl(C(i),\, 0.5\bigr)}{C(0)}
  \label{eq:minetti_factor}
\end{equation}
where the floor at $0.5$ prevents unrealistic negative costs on steep descents.

During anisotropic Dijkstra pathfinding, the directional slope between adjacent cells $(r_1, c_1) \to (r_2, c_2)$ is computed from the DEM-derived slope gradients.
At each grid cell, the partial derivatives $\partial z / \partial x$ and $\partial z / \partial y$ are computed via central finite differences on the DEM.
For a step between two adjacent cells, these derivatives are averaged to obtain a representative gradient at the midpoint:
\begin{equation}
  \bar{g}_x = \tfrac{1}{2}\!\left(\frac{\partial z}{\partial x}\bigg|_{(r_1,c_1)} + \frac{\partial z}{\partial x}\bigg|_{(r_2,c_2)}\right), \quad
  \bar{g}_y = \tfrac{1}{2}\!\left(\frac{\partial z}{\partial y}\bigg|_{(r_1,c_1)} + \frac{\partial z}{\partial y}\bigg|_{(r_2,c_2)}\right)
  \label{eq:gradavg}
\end{equation}
Let $\Delta c = c_2 - c_1$ and $\Delta r = r_2 - r_1$ be the step direction.
The signed slope gradient along the direction of travel is:
\begin{equation}
  i_{12} = \frac{\bar{g}_x \cdot \Delta c + \bar{g}_y \cdot \Delta r}{\sqrt{\Delta c^2 + \Delta r^2}}
  \label{eq:dirslope}
\end{equation}
This projects the terrain gradient onto the movement vector, yielding a positive value for uphill travel and negative for downhill.
The traversal cost of a single Dijkstra step from cell $(r_1, c_1)$ to an adjacent cell $(r_2, c_2)$ is:
\begin{equation}
  d_{12} = \|\Delta\| \cdot \frac{\text{Cost}_{r_1 c_1} + \text{Cost}_{r_2 c_2}}{2} \cdot \phi(i_{12})
  \label{eq:edgecost}
\end{equation}
where $\text{Cost}_{rc}$ is the combined base cost at cell $(r,c)$ from \cref{eq:combined}, $\phi(i_{12})$ is the Minetti slope factor from \cref{eq:minetti_factor}, and $\|\Delta\| = \sqrt{\Delta c^2 + \Delta r^2}$ is the Euclidean step length between adjacent cells ($\|\Delta\| = 1$ for cardinal moves, $\|\Delta\| = \sqrt{2}$ for diagonal moves in the 8-connected grid).
The base cost is averaged between the source and destination cells to smooth transitions across terrain boundaries.

This produces an asymmetric traversal cost: the step cost $d_{12}$ from cell~1 to cell~2 differs from $d_{21}$ in the reverse direction, because the Minetti factor satisfies $\phi(i) \neq \phi(-i)$ for any non-zero slope.

% ── 2.5 Cost Matrix (moved before fatigue so C_ij is defined first) ──
\subsection{Asymmetric Cost Matrix Construction}
\label{sec:costmatrix}

Given $n$ control points plus the depot (node~0), the pairwise cost matrix $\mathbf{C} \in \mathbb{R}^{(n+1) \times (n+1)}$ is computed by running anisotropic Dijkstra from each node on the combined cost grid (\cref{eq:combined}) with Minetti slope factors (\cref{eq:minetti_factor}).
The per-step traversal cost used as the Dijkstra edge weight is given by \cref{eq:edgecost}, which integrates the base cost raster, the directional Minetti multiplier, and the step geometry into a single directional cost per grid move.

For each node $i \in V$, a single-source anisotropic Dijkstra is run on the (optionally downsampled) grid, starting from the grid cell corresponding to node $i$.
The algorithm explores all 8-connected neighbours using a priority queue, computing the traversal cost to each adjacent cell via \cref{eq:edgecost}.
Because each step cost is directional (as described above), the shortest-path cost from node $i$ to node $j$ generally differs from the cost from $j$ to $i$, yielding an asymmetric cost matrix.

The cost matrix entry $C_{ij}$ is then the shortest-path distance from node $i$'s grid cell to node $j$'s grid cell in this directed cost field.
Running $n+1$ independent Dijkstra computations (one per node) yields the full matrix.
Cells with base cost $\geq 9.0$ (impassable features) are treated as hard barriers with a fixed penalty cost, preventing routes from crossing cliffs, water, or out-of-bounds areas.

The grid is optionally downsampled by factor $s$, averaging the base cost in $s \times s$ blocks and rescaling node coordinates accordingly.
This reduces computation time from $O(nHW \log(HW))$ to $O\bigl(n \frac{HW}{s^2} \log \frac{HW}{s^2}\bigr)$ at the cost of path resolution.

The resulting matrix satisfies $C_{ij} \neq C_{ji}$ in general, with asymmetry measured as:
\begin{equation}
  \mathcal{A} = \frac{\text{mean}_{i \neq j}\, |C_{ij} - C_{ji}|}{\text{mean}_{i \neq j}\, C_{ij}} \times 100\%
  \label{eq:asymmetry}
\end{equation}

% ── 2.6 Fatigue ──
\subsection{Cumulative Fatigue Model}
\label{sec:fatigue}

\paragraph{Revision note.}
An earlier version of this model defined fatigue as $f(t) = 1 + \lambda\, t/B$, a multiplier driven purely by cumulative \emph{base cost} $t$ (i.e.\ elapsed effort, in the same units as $B$).
That model is provably order-invariant: expanding the total fatigue-adjusted cost of a route algebraically collapses to $S + (\lambda/2B)(S^2 - Q)$, where $S = \sum_\ell C_\ell$ and $Q = \sum_\ell C_\ell^2$ over the route's legs---a function of the \emph{multiset} of arcs used, independent of visit order.
Consequently the AOPF as originally posed reduces to a classical Orienteering Problem with a nonlinear (quadratic) budget, and none of the arrival-time machinery in \cref{sec:formulation} (McCormick linearisation, flow-based time tracking, position-dependent cuts) is actually necessary to solve it.
The model below replaces the elapsed-cost multiplier with a per-node \emph{fatigue state} driven by the terrain actually traversed---elevation gain, elevation loss, and distance---which breaks this collapse (\cref{sec:order-dependence}).
\Cref{sec:formulation} (including its cutting-plane subsections) still describes the arrival-time formulation built for the original (collapsing) model; adapting it to the state-based model below is future work (\cref{sec:discussion}).

\paragraph{Fatigue state.}
Rather than a single scalar elapsed-cost fraction, we track a per-node fatigue state, accumulated along the path actually taken to reach $i$.
Each arc $(i,j)$ contributes a signed increment:
\begin{equation}
  \psi_{ij} = \varphi^+_{ij} - \rho\, \varphi^-_{ij} + \mu\, \delta_{ij}
  \label{eq:psi}
\end{equation}
where $\varphi^+_{ij}$ and $\varphi^-_{ij}$ are the cumulative uphill and downhill elevation change (metres) along the cost-optimal path for arc $(i,j)$ (the same path underlying $C_{ij}$, \cref{sec:costmatrix}), and $\delta_{ij}$ is the path length (metres) along that same path.
$\rho \in [0,1]$ is the fraction of downhill elevation that offsets accumulated fatigue---a physiological unknown (the trade-off between metabolic recovery and eccentric muscular damage on descents is not settled in the literature we are aware of) and is therefore treated as a sensitivity parameter, swept rather than fixed to a single value (\cref{sec:experiments}).
$\mu \geq 0$ converts horizontal distance into vertical-metre-equivalent fatigue units; unlike $\rho$ it is not a free parameter but is derived from the Minetti curve already underlying $C_{ij}$ (\cref{eq:minetti_factor}):
\begin{equation}
  \mu = \dfrac{c(0)}{\min_{i > 0}\; c(i)/i}
  \label{eq:mu}
\end{equation}
i.e.\ the metabolic cost of one flat metre divided by the metabolic cost of the cheapest-to-reach vertical metre (the gradient $i^\star$ that minimises cost per unit of \emph{vertical} rise).
Both $\varphi^+_{ij}/\varphi^-_{ij}$ and $\delta_{ij}$ are accumulated as byproducts of the same anisotropic Dijkstra used to build $C_{ij}$, at negligible extra cost.

The true fatigue state reflects a physiological floor: fatigue cannot recover below zero.
Writing $F_0 = 0$ at the depot,
\begin{equation}
  F_j = \max\bigl(0,\, F_i + \psi_{ij}\bigr)
  \label{eq:fatigue_clip}
\end{equation}
for the arc $(i,j)$ actually taken to reach $j$.
The effective (fatigue-adjusted) cost of arc $(i,j)$ departing node $i$ at state $F_i$ is:
\begin{equation}
  C^f_{ij}(F_i) = C_{ij} \cdot \bigl(1 + \lambda\, F_i\bigr)
  \label{eq:fatigue_cost}
\end{equation}
and the total fatigue-adjusted cost of a route $\sigma = (0, v_1, \ldots, v_k, 0)$ is $D^f(\sigma) = \sum_{\ell} C^f_{v_\ell, v_{\ell+1}}(F_{v_\ell})$, with $F$ propagated leg-by-leg via \cref{eq:fatigue_clip}.
Because $\lambda$ now multiplies a state expressed in metres of elevation/distance rather than a $[0,1]$-normalised elapsed-budget fraction, it is not dimensionally comparable to the $\lambda$ of the original model and must be recalibrated (\cref{sec:experiments}).

\paragraph{Order-dependence.}
\label{sec:order-dependence}
Unlike the original model, $D^f(\sigma)$ is not invariant to the order in which a fixed set of arcs is traversed.
Expanding $D^f$ (ignoring the clip in \cref{eq:fatigue_clip} for the moment, i.e.\ using the unclipped state $G_i = \sum_{m < \ell} \psi_m$) gives a cross term $\lambda \sum_{m < \ell} C_\ell\, \psi_m$.
In the original model this collapsed because the accumulator was the same quantity being scaled ($\psi_m \propto C_m$ for every arc, making the term $\lambda \sum_{m<\ell} C_\ell C_m$, symmetric under any permutation).
Here $\psi_{ij}$ (elevation- and distance-driven) is not proportional to $C_{ij}$ (which also depends on the areal terrain weight, \cref{eq:acr}) for arbitrary arcs, so the cross term is generically asymmetric under reordering: two routes visiting the same nodes via the same arcs in different sequences can have substantially different $D^f(\sigma)$.
The $\max(0, \cdot)$ clip in \cref{eq:fatigue_clip} is a second, independent source of order-dependence on top of this (a downhill-heavy prefix can floor the state in a way a later downhill leg cannot undo retroactively), but---as verified numerically on the implementation---order-dependence already holds without it, as long as $\psi_{ij}/C_{ij}$ is not constant across arcs.
This is why the distance term in \cref{eq:psi} is defined using raw path length $\delta_{ij}$ rather than a cost-derived quantity: a $\mu$ term correlated with $C_{ij}$ would reintroduce (a component of) the collapsing cross term, diluting the very property this model exists to establish.


% ════════════════════════════════════════════════════════════════
\section{Mathematical Formulation}
\label{sec:formulation}
% ════════════════════════════════════════════════════════════════

We formulate the Asymmetric Orienteering Problem with Fatigue (AOPF) as a mixed-integer linear program.
The core challenge is the bilinear fatigue term $G_i \cdot x_{ij}$ in the budget constraint, where $G_i$ is the (unclipped) fatigue state at node $i$ (\cref{sec:fatigue}).
We present two formulations that handle this term differently: an MTZ-based formulation that linearises it via McCormick envelopes, and a flow-based formulation where the fatigue budget is natively linear.
Both share the same objective, arc/node variables, and subtour elimination cuts; they differ only in the state-tracking variables and the fatigue budget constraint.

\subsection{Sets and Parameters}

Let $V = \{0, 1, \ldots, n\}$ be the set of nodes, where node~0 is the depot.
The parameters $C_{ij}$, $\psi_{ij}$, $B$, $\lambda$, $\rho$, and $\mu$ are as defined in \cref{sec:cost,sec:fatigue}.
Each node $i$ has a score $p_i \geq 0$ (with $p_0 = 0$).

\subsection{Decision Variables}

Both formulations share arc variables $x_{ij} \in \{0,1\}$ and node-visit variables $y_i \in \{0,1\}$.
They differ in how the fatigue state and budget are modelled.

\subsection{Objective}

Maximise total score:
\begin{equation}
  \max \sum_{i \in V} p_i \, y_i
  \label{eq:obj}
\end{equation}

\subsection{Shared Constraints}

\paragraph{Flow conservation.}
Each visited node has exactly one incoming and one outgoing arc:
\begin{align}
  \sum_{j \in V \setminus \{i\}} x_{ji} &= y_i && \forall\, i \in V \label{eq:inflow} \\
  \sum_{j \in V \setminus \{i\}} x_{ij} &= y_i && \forall\, i \in V \label{eq:outflow}
\end{align}

\paragraph{Depot.}
The depot is always visited: $y_0 = 1$.

\paragraph{Arc elimination.}
Arcs that are structurally infeasible are fixed to zero a priori.
An arc $(i,j)$ is eliminated if the base cost alone exceeds the budget:
\begin{equation}
  x_{ij} = 0 \quad \text{if } C_{ij} + C_{j0} > B \;\text{ or }\; C_{0i} + C_{ij} + C_{j0} > B
  \label{eq:arcelim}
\end{equation}
The original (linear) model additionally applied a fatigue-aware pre-filter, testing whether the direct route $0 \to i \to j \to 0$ remained budget-feasible once fatigue was accounted for.
Under the revised model (\cref{sec:fatigue}) this test is not available as a cheap a~priori filter: evaluating it requires an upper bound $\hat{G}_i$ on the fatigue state reachable at $i$, but $\hat{G}_i$ is itself only well-defined once the surviving-arc graph is fixed (\cref{eq:ghat} below)---testing arc $(i,j)$ against a bound that depends on which arcs survive, including $(i,j)$ itself, is circular.
We therefore retain only the base-cost filter~\eqref{eq:arcelim}; it is a conservative superset of the true survivors (some arcs it keeps may still be fatigue-infeasible), which is exactly what \cref{eq:ghat} needs to remain a valid bound.

\paragraph{Fatigue-state bounds $\hat{G}_i$ and $\check{G}_i$.}
Both formulations below need a valid upper bound $\hat{G}_i$ on the (unclipped) fatigue state reachable at $i$, to size big-$M$ constants and McCormick/coupling bounds---the role $B - C_{0i} - C_{j0}$ played for arrival times in the original model.
Unlike elapsed cost, $\psi_{ij}$ can be negative, so the surviving-arc graph (\cref{eq:arcelim}) weighted by $\psi$ can contain positive-weight cycles, making the longest-walk problem that would naturally define $\hat{G}_i$ NP-hard in general.
We instead compute $\hat{G}_i$ by a Bellman--Ford relaxation in the $(\max, +)$ semiring, capped at $n-1$ rounds:
\begin{equation}
  \hat{G}_i = \max\Bigl(0,\ \max_{\substack{\text{walks } 0 \to i \\ \text{of length} \leq n-1}} \sum_{\text{legs}} \psi\Bigr)
  \label{eq:ghat}
\end{equation}
Any \emph{simple} path through the $n+1$ nodes has at most $n-1$ arcs, so capping the relaxation at $n-1$ rounds is a superset of all simple-path sums---the bound remains valid regardless of cycle sign, and the relaxation is well-defined (no divergence) despite positive cycles, since we simply stop before they can be exploited.
An analogous $\check{G}_i$, a valid \emph{lower} bound needed only by Formulation~B (below), is computed the same way with a $(\min,+)$ relaxation.
Both run in $O(n^3)$ (\,$O(n^2)$ arcs $\times$ $O(n)$ rounds), trivial for the instance sizes here.
$\hat{G}_i$ and $\check{G}_i$ are constructed to be valid bounds for the LP relaxation's own state variables (defined below); the route an accepted solution actually reports is always re-checked directly against the exact clipped model~\eqref{eq:fatigue_clip} before being output; see \cref{sec:discussion} for further remarks on the relationship between the two.

\subsection{Formulation A: MTZ with McCormick Linearisation}
\label{sec:formulation_mtz}

This formulation introduces fatigue-state variables $G_i \in [0, \hat{G}_i]$ at each node and auxiliary variables $u_{ij} \geq 0$ to linearise the bilinear fatigue term $G_i \cdot x_{ij}$.

\paragraph{State propagation (MTZ-type).}
For each arc $(i,j)$ with $j \neq 0$:
\begin{equation}
  G_j \geq G_i + \psi_{ij} - M_{ij}(1 - x_{ij})
  \label{eq:mtz}
\end{equation}
where the big-$M$ constant is:
\begin{equation}
  M_{ij} = \hat{G}_i + \max(\psi_{ij},\, 0)
  \label{eq:bigM}
\end{equation}
so that the constraint is vacuous when $x_{ij}=0$ (the largest $G_i$ can be is $\hat{G}_i$, and the largest $\psi_{ij}$ can contribute is $\max(\psi_{ij},0)$).
Because $G_i$ additionally has a hard column floor of $0$ and~\eqref{eq:mtz} is a one-sided inequality (not an equality), the LP is free to set $G_j$ to the tightest value satisfying every incoming constraint---which, combined with the budget constraint's pressure to keep $G$ small, reproduces the $\max(0, \cdot)$ clip of~\eqref{eq:fatigue_clip} along the path actually selected by $x$, without needing to encode the clip explicitly.

\paragraph{McCormick linearisation.}
The fatigue budget involves the bilinear term $G_i \cdot x_{ij}$.
We introduce $u_{ij} \approx G_i \cdot x_{ij}$ with the standard McCormick envelope, using $\hat{G}_i$ as the bound:
\begin{align}
  u_{ij} &\geq G_i - \hat{G}_i(1 - x_{ij}) \label{eq:mc1} \\
  u_{ij} &\leq \hat{G}_i \, x_{ij} \label{eq:mc2} \\
  u_{ij} &\leq G_i \label{eq:mc3}
\end{align}
Constraint~\eqref{eq:mc1} forces $u_{ij} = G_i$ when $x_{ij} = 1$; constraint~\eqref{eq:mc2} forces $u_{ij} = 0$ when $x_{ij} = 0$; and constraint~\eqref{eq:mc3} ensures $u_{ij} \leq G_i$.

\paragraph{Fatigue budget (MTZ).}
Unlike the original model, there is no $B$ to normalise by: $G_i$ is expressed in metres, not as an elapsed-budget fraction.
\begin{equation}
  \sum_{(i,j) \in A} C_{ij} \, x_{ij} + \lambda \sum_{(i,j) \in A} C_{ij} \, u_{ij} \leq B
  \label{eq:fatigue_budget}
\end{equation}
The McCormick envelope is an outer approximation of the true bilinear set, introducing LP relaxation looseness that grows with the fatigue rate $\lambda$, as before.

\subsection{Formulation B: Single-Commodity Flow}
\label{sec:formulation_flow}

This formulation replaces the MTZ state variables and McCormick auxiliary variables with flow variables $g_{ij}$.
We define $g_{ij}$ so that, by construction of the constraints below, $g_{ij} = G_i \cdot x_{ij}$ at any integer-feasible solution.

\paragraph{Flow-arc coupling.}
Unlike the original model, where elapsed cost is always non-negative, $\psi_{ij}$ can be negative (net downhill recovery), so $g_{ij}$ needs \emph{both} an upper and a lower coupling bound:
\begin{equation}
  \check{G}_i \cdot x_{ij} \leq g_{ij} \leq \hat{G}_i \cdot x_{ij} \qquad \forall\, (i,j) \in A
  \label{eq:flow_coupling}
\end{equation}
A single-sided bound $g_{ij} \leq \hat{G}_i x_{ij}$ (mirroring the original model's~\eqref{eq:flow_coupling}) would be unsound here: combined with the conservation equality below, it can force $g_{ij}$ to be \emph{inflated} on arcs upstream of a strongly-recovering leg, simply to keep the (non-negative-only) flow variable feasible against a negative increment further down the path---an artifact of the equality system, not a reflection of the true state.
The two-sided bound~\eqref{eq:flow_coupling} gives $g_{ij}$ room to actually hold a negative increment.

\paragraph{Fatigue flow conservation.}
At each non-depot node $j$:
\begin{equation}
  \sum_{i \in V} \bigl(g_{ij} + \psi_{ij} \, x_{ij}\bigr) = \sum_{k \in V} g_{jk} \qquad \forall\, j \in V \setminus \{0\}
  \label{eq:flow_conservation}
\end{equation}
Depot outgoing flow is fixed to zero: $g_{0j} = 0$ for all $j$.
This has the same structure as the original model's time-flow conservation, but the coupling bounds~\eqref{eq:flow_coupling} being two-sided (rather than $g_{ij} \geq 0$) is what makes it well-posed for a $\psi$ that changes sign.
The flow structure eliminates subtours for integer solutions; the explicit SECs and connectivity cuts (\cref{eq:sec,eq:outcut,eq:incut}) remain necessary to cut off fractional subtours in the LP relaxation.

\paragraph{Fatigue budget (flow).}
\begin{equation}
  \sum_{(i,j) \in A} C_{ij} \, x_{ij} + \lambda \sum_{(i,j) \in A} C_{ij} \, g_{ij} \leq B
  \label{eq:fatigue_budget_flow}
\end{equation}
As in the original model, the flow formulation has $O(n^2)$ additional continuous variables (one $g_{ij}$ per arc) compared to $O(n)$ state variables in the MTZ formulation, but produces a tighter LP relaxation than the McCormick envelope.

\subsection{Subtour Elimination and Connectivity Cuts}

Both formulations use the same family of cutting planes, added dynamically during the B\&C search.
Because the cost matrix is asymmetric ($C_{ij} \neq C_{ji}$), all cuts are formulated on the directed arc variables $x_{ij}$.
An inequality is \emph{valid} for the AOPF if it is satisfied by every integer-feasible solution; adding valid inequalities to the LP relaxation tightens the upper bound without excluding any feasible solution.

\paragraph{Subtour elimination constraints (SECs).}
For a subset $S \subseteq V \setminus \{0\}$ (identified by the separation routines in \cref{sec:separation}), the directed SEC prevents cycles disconnected from the depot:
\begin{equation}
  \sum_{i \in S} \sum_{j \in S} x_{ij} \leq |S| - 1
  \label{eq:sec}
\end{equation}

\paragraph{Connectivity cuts.}
These ensure that every visited node can reach and be reached from the depot:
\begin{align}
  \sum_{i \in S} \sum_{j \notin S} x_{ij} &\geq y_k && \forall\, k \in S \label{eq:outcut} \\
  \sum_{i \notin S} \sum_{j \in S} x_{ij} &\geq y_k && \forall\, k \in S \label{eq:incut}
\end{align}
A combined two-way cut further tightens the relaxation:
\begin{equation}
  \sum_{i \in S} \sum_{j \notin S} x_{ij} + \sum_{i \notin S} \sum_{j \in S} x_{ij} \geq 2\, y_k \quad \text{for some } k \in S
  \label{eq:twocut}
\end{equation}

\subsection{Budget-Based Cutting Planes}

The cuts in this subsection exploit the budget constraint to eliminate LP solutions that are structurally infeasible due to excessive routing cost.
They complement the subtour and connectivity cuts above, which enforce tour structure but do not account for cost.
We present four families of budget-based cuts, numbered (B1)--(B4).
The separation routines used to identify violated instances of each cut are described in \cref{sec:separation}.

\paragraph{(B1) Lifted cover cuts.}
The fatigue budget constraint (\cref{eq:fatigue_budget,eq:fatigue_budget_flow}) limits the total cost of arcs used in the solution.
Consider a subset of arcs $\mathcal{C} \subseteq A$ whose combined base costs exceed the budget: $\sum_{(i,j) \in \mathcal{C}} C_{ij} > B$.
A feasible solution may use some arcs from $\mathcal{C}$, but it cannot use all of them simultaneously---doing so would exceed the budget even before accounting for fatigue.
Such a set $\mathcal{C}$ is called a \emph{cover}, and the corresponding \emph{cover inequality} states:
\begin{equation}
  \sum_{(i,j) \in \mathcal{C}} x_{ij} \leq |\mathcal{C}| - 1
  \label{eq:cover}
\end{equation}
That is, at least one arc in the cover must be excluded from any feasible solution.
Each cover can be \emph{lifted}: for arcs $e \notin \mathcal{C}$, a coefficient $\alpha_e \geq 0$ is computed such that adding $\alpha_e\, x_e$ to the left-hand side preserves validity, strengthening the inequality.

\paragraph{(B2) Routing infeasibility cuts.}
The cover cuts (B1) operate on arc variables $x_{ij}$, which the LP can soften through fractional values.
We introduce a complementary family of inequalities defined purely on the node-visit variables $y_i$, exploiting the fact that certain subsets of nodes cannot all be visited within the budget regardless of the route chosen.

For a subset $Q \subseteq V \setminus \{0\}$, let $C_{\text{assign}}(Q)$ denote the optimal value of the assignment relaxation on the subgraph induced by $Q \cup \{0\}$: the minimum total base cost of a set of arcs where every node has exactly one incoming and one outgoing arc, solvable in $O(|Q|^3)$ by the Hungarian algorithm.
Since fatigue only increases travel cost, $C_{\text{assign}}(Q)$ is a lower bound on the fatigue-adjusted cost of any cycle through $Q \cup \{0\}$.

\begin{proposition}[Routing infeasibility]
\label{prop:routing}
For any $Q \subseteq V \setminus \{0\}$ with $C_{\text{assign}}(Q) > B$, the inequality
\begin{equation}
  \sum_{i \in Q} y_i \leq |Q| - 1
  \label{eq:routing_infeasibility}
\end{equation}
is satisfied by every integer-feasible solution of the AOPF\@.
\end{proposition}

\begin{proof}
Suppose an integer-feasible solution has $y_i = 1$ for all $i \in Q$.
The solution defines a cycle $\sigma$ through the depot visiting a set $W$ with $Q \subseteq W$.
By definition, $D^f(\sigma) \geq C_{\text{assign}}(Q)$ since $\sigma$ visits all nodes in $Q \cup \{0\}$ and $C_{\text{assign}}(Q)$ is a lower bound on the base cost of any such cycle, and fatigue only increases cost.
But $D^f(\sigma) \leq B$ (budget feasibility), contradicting $C_{\text{assign}}(Q) > B$.
\end{proof}

Unlike cover cuts (B1), which the LP can satisfy by reducing fractional arc values, routing infeasibility cuts constrain the $y_i$ variables directly.
The LP cannot circumvent them by redistributing arc flow, since the inequality involves no arc variables.

\paragraph{Remark.}
Propositions~\ref{prop:routing} and~\ref{prop:cyclecover} (below) are stated purely in terms of $D^f(\sigma)$ and hold for \emph{any} route-cost function satisfying $D^f(\sigma) \geq \sum_{(i,j) \in \sigma} C_{ij}$ (fatigue never reduces cost)---true both of the original model and of the revised one (\cref{sec:fatigue}), since $C^f_{ij}(F_i) = C_{ij}(1 + \lambda F_i) \geq C_{ij}$ as $F_i \geq 0$ always.
Their proofs therefore carry over unchanged; only the definition of $D^f$ itself has changed.

\paragraph{(B3) Directed cycle cover cuts.}
We adapt the cycle cover inequalities of \citet{fischetti1998solving} from the undirected OP to our directed formulation.
Let $F \subset A$ be a set of directed arcs forming a directed cycle through nodes $V(F)$ whose base cost $\sum_{(i,j) \in F} C_{ij}$ exceeds $B$.

\begin{proposition}[Directed cycle cover]
\label{prop:cyclecover}
For any directed cycle $F \subset A$ with $\sum_{(i,j) \in F} C_{ij} > B$, the inequality
\begin{equation}
  \sum_{(i,j) \in F} x_{ij} \leq \sum_{v \in V(F)} y_v - 1
  \label{eq:cycle_cover}
\end{equation}
is satisfied by every integer-feasible solution of the AOPF\@.
\end{proposition}

\begin{proof}
Suppose for contradiction that an integer-feasible solution has $\sum_{(i,j) \in F} x_{ij} = \sum_{v \in V(F)} y_v$, i.e., every arc in $F$ is used and every node in $V(F)$ is visited.
Then the route contains all arcs of $F$, whose base cost alone exceeds $B$.
Since arc costs are non-negative, $D^f(\sigma) \geq \sum_{(i,j) \in F} C_{ij} > B$, contradicting budget feasibility.
Therefore $\sum_{(i,j) \in F} x_{ij} \leq \sum_{v \in V(F)} y_v - 1$ in every integer-feasible solution.
\end{proof}

Inequality~\eqref{eq:cycle_cover} is strictly stronger than the cover (B1) because the right-hand side involves the $y_v$ variables: when the LP sets $y_v < 1$ for some $v \in V(F)$, the right-hand side decreases, forcing even fewer arcs to be used.
Cover cuts (B1) remain useful because their separation is computationally cheaper.

\paragraph{(B4) Directed path inequality cuts.}
We adapt the path inequalities of \citet{fischetti1998solving} to the directed asymmetric OP with fatigue.
Let $P = (v_1 \to v_2 \to \cdots \to v_k)$ be a directed path through $k$ non-depot nodes, where each $v_\ell \in V \setminus \{0\}$ is a node and the arcs $(v_\ell, v_{\ell+1})$ for $\ell = 1, \ldots, k-1$ are traversed in sequence.
Let $V(P) = \{v_1, \ldots, v_k\}$ denote the node set of $P$.
Propagate the exact (clipped) fatigue state deterministically along $P$ using the base matrices: $\phi_1 = \max(0, \psi_{0,v_1})$ and $\phi_\ell = \max(0,\, \phi_{\ell-1} + \psi_{v_{\ell-1}, v_\ell})$ for $\ell = 2, \ldots, k$, matching~\eqref{eq:fatigue_clip} exactly since $P$ is a single fixed sequence with no branching.
Let $\Phi(P) = C_{0,v_1} \cdot (1 + \lambda \cdot 0) + \sum_{\ell=1}^{k-1} C_{v_\ell, v_{\ell+1}}^f(\phi_\ell)$ be the fatigue-adjusted cost of traversing $P$ from the depot (the state is $0$ on departure, so the first leg is unpenalised).
Define the \emph{feasible closing set}:
\begin{equation}
  W(P) = \left\{ v \in V \setminus V(P) : \Phi(P) + C_{v_k, v}^f(\phi_k) + C_{v, 0}^f\bigl(\max(0, \phi_k + \psi_{v_k, v})\bigr) \leq B \right\}
  \label{eq:wp}
\end{equation}
where $C_{i,j}^f(\cdot)$ is the fatigue-adjusted cost of arc $(i,j)$ as defined in~\eqref{eq:fatigue_cost}, evaluated at the deterministic state reached at each point along $P$ then closing through $v$.
The set $W(P)$ contains exactly those nodes through which the path can be completed into a budget-feasible cycle.
The path inequality states:
\begin{equation}
  \sum_{\ell=1}^{k-1} x_{v_\ell, v_{\ell+1}} - \sum_{v \in V(P)} y_v + y_{v_1} + y_{v_k} - \sum_{v \in W(P)} x_{v_k, v} \leq 0
  \label{eq:path_ineq}
\end{equation}
Intuitively, if all path arcs are used (first sum is large), then either some interior path nodes are unvisited (second and third terms compensate), or the endpoint $v_k$ connects to a node in $W(P)$ (last sum compensates).
If neither holds, the path cannot be completed within budget.
As in the original model, the fatigue-inflated return cost from $v_k$ restricts $W(P)$ relative to the fatigue-free case, producing tighter cuts on high-fatigue instances; the effect is now asymmetric in $\rho$ as well as in $\lambda$, since a downhill-heavy prefix leaves more headroom in $\phi_k$ than an uphill-heavy one of the same base cost.

\subsection{Separation Routines}
\label{sec:separation}

The valid inequalities defined in the preceding subsections are added dynamically during the B\&C search.
At each node of the search tree, the solver iterates through a cut loop that attempts to find violated inequalities in the current LP solution $(y^*, x^*, f^*)$.
The loop terminates when no violated cuts are found or a maximum number of iterations is reached.
Below we describe how violated instances of each cut family are identified.

\paragraph{Subtour and connectivity cuts.}
A directed support graph $H^* = (V^*, A^*)$ is constructed from arcs with $x_{ij}^* > 0.5$.
Two complementary procedures are applied:
\begin{enumerate}[nosep]
  \item \emph{Kosaraju SCC detection.}
        Kosaraju's algorithm identifies all strongly connected components of $H^*$ in $O(|V^*| + |A^*|)$ time.
        Any SCC not containing the depot yields a violated subset $S$, and all four cut types~(\cref{eq:sec,eq:outcut,eq:incut,eq:twocut}) are added for $S$.
  \item \emph{Reverse BFS.}
        A reverse BFS from the depot identifies nodes that can reach the depot in $H^*$.
        Any visited node ($y_i^* > 0.5$) not in this reachable set requires an outgoing connectivity cut.
\end{enumerate}
This two-phase separation is essential for the asymmetric OP\@: Kosaraju detects closed subtours, while the reverse BFS catches paths that leave the depot but cannot return.

\paragraph{Max-flow connectivity cuts.}
The rounding-based separation above (using threshold $x_{ij}^* > 0.5$) may miss violated connectivity cuts when the fractional arc values are spread thinly.
To strengthen the separation, for each visited node $k$ with $y_k^* > 0.1$, a max-flow computation (Dinic's algorithm) is performed from the depot to $k$ on the fractional support graph, using $x_{ij}^*$ as arc capacities.
If the max-flow value is less than $y_k^*$, the min-cut identifies a subset $S$ with $k \in S$ and $0 \notin S$ for which the connectivity cut~\eqref{eq:outcut} is violated.
This exact separation complements the heuristic rounding-based approach and is particularly effective on asymmetric instances where fractional flow is distributed across many arcs.

\paragraph{(B1) Cover cuts.}
Arcs with $x_{ij}^* > 0$ are sorted by decreasing $x_{ij}^*$ and accumulated greedily until their total base cost exceeds $B$.
The resulting set $\mathcal{C}$ defines a violated cover if $\sum_{(i,j) \in \mathcal{C}} x_{ij}^* > |\mathcal{C}| - 1$, i.e., the LP uses more of these arcs than any feasible solution allows.
Each violated cover is optionally strengthened via sequential lifting.
Up to 3 cover cuts are added per cut-loop iteration; adding more would slow the LP solve without proportional improvement in the bound.
To account for fatigue, the arc weights used in the greedy accumulation are inflated from the base cost $C_{ij}$ to $C_{ij}\bigl(1 + \lambda\, \max(\check{G}_i, 0)\bigr)$: the true clipped state satisfies $F_i \geq \check{G}_i$ always (as $\check{G}_i$ is a valid lower bound on the unclipped state, \cref{eq:ghat}), and $F_i \geq 0$ trivially, so $\max(\check{G}_i,0)$ is the tightest simple lower bound on the minimum possible fatigue state at $i$ available without solving for the specific route.
This produces smaller covers---fewer arcs are needed to exceed $B$---yielding tighter inequalities, and is only useful (multiplier $>1$) on arcs where $\check{G}_i$ is positive, i.e.\ where every budget-admissible route reaching $i$ is net uphill.
The original model instead inflated weights by $C_{ij}(1 + \lambda\, C_{0i}/B)$, using the direct depot cost $C_{0i}$ as a proxy for minimum elapsed time; that proxy has no valid meaning under the revised, elevation-driven model and is not used here.

\paragraph{(B2) Routing infeasibility cuts.}
Nodes with $y_i^* > 0.5$ are sorted by decreasing $y_i^*$ and added one by one to a candidate set $Q$.
After each addition, the assignment bound $C_{\text{assign}}(Q)$ is computed via the Hungarian algorithm ($O(|Q|^3)$).
When $C_{\text{assign}}(Q) > B$ and $\sum_{i \in Q} y_i^* > |Q| - 1$, a violated routing infeasibility cut has been found.
The set $Q$ is then minimised: nodes are removed in order of increasing $y_i^*$ while $C_{\text{assign}}(Q)$ remains above $B$, yielding a minimal infeasible set and hence the tightest cut.

\paragraph{(B3) Cycle cover cuts.}
Short directed cycles (length 3--4) are enumerated in the support graph $H^*$.
For each cycle $F$, the base cost $\sum_{(i,j) \in F} C_{ij}$ is computed.
If $\sum_{(i,j) \in F} C_{ij} > B$ and the inequality~\eqref{eq:cycle_cover} is violated by $(x^*, y^*)$, the cut is added.

\paragraph{(B4) Path inequality cuts.}
Directed paths of length 2--3 are enumerated in $H^*$.
For each path $P = (v_1 \to \cdots \to v_k)$, the feasible closing set $W(P)$ is computed using the base matrices $C_{ij}$, $\psi_{ij}$ (not the LP solution values): the fatigue state at each path node is accumulated deterministically via the clip recursion $\phi_\ell = \max(0, \phi_{\ell-1} + \psi_{v_{\ell-1},v_\ell})$, and the fatigue-adjusted cost of each subsequent arc is evaluated at this state, as in~\eqref{eq:wp}.
If $\Phi(P)$ (the path's own fatigue-adjusted cost) already exceeds $B$, $W(P)$ is empty, or a direct return $v_k \to 0$ is itself budget-feasible (in which case the path inequality would add nothing, since closing directly is always available), the path is skipped.
Otherwise the cut is added if inequality~\eqref{eq:path_ineq} is violated by the current LP solution.

% ════════════════════════════════════════════════════════════════
\section{Solution Methods}
\label{sec:methods}
% ════════════════════════════════════════════════════════════════

\subsection{Simulated Annealing (SA)}
\label{sec:sa}

\paragraph{Method selection.}
The SA metaheuristic was selected after a comparative study on 21 benchmark instances across varying sizes ($n = 20$--$100$), asymmetry levels, and fatigue rates.
Four methods were evaluated with equal computational budget (2 seconds per method per instance, 10 independent seeds for stochastic methods): Greedy Insertion (deterministic), Genetic Algorithm (GA, population 50), Ant Colony Optimisation (ACO, 30 ants), and Simulated Annealing (SA).

\begin{table}[H]
\centering
\caption{Heuristic comparison: mean gap below SA (\%) across 21 benchmark instances (equal 2\,s time budget, 10 seeds per stochastic method). All differences are statistically significant (Wilcoxon signed-rank, $p < 0.001$).}
\label{tab:heuristic_comparison}
\begin{tabular}{lrrrr}
\toprule
\textbf{Method} & \textbf{Mean gap} & \textbf{Std} & \textbf{Min gap} & \textbf{Max gap} \\
\midrule
Greedy  & 23.0\% &  9.4\% &  3.7\% & 40.0\% \\
GA      & 49.2\% & 11.0\% & 24.6\% & 71.9\% \\
ACO     &  8.1\% &  4.9\% &  1.3\% & 24.1\% \\
SA      &  ---   & ---    &  ---   & ---    \\
\bottomrule
\end{tabular}
\end{table}

SA significantly outperforms all three alternatives under equal time budgets (Wilcoxon signed-rank test, $n = 21$ instances, all $p < 0.001$; SA achieves a higher mean score than each competitor on all 21 instances).
The GA performs worst (49.2\% below SA on average), as crossover on directed routes frequently produces infeasible offspring requiring expensive repair, and the rugged fitness landscape under asymmetric costs degrades convergence.
ACO is the closest competitor (8.1\% below SA) but its pheromone mechanism weakens on asymmetric instances where edge quality is context-dependent on elapsed time and travel direction---the gap widens to 24.1\% on the hardest instance.
Greedy Insertion (23.0\% below SA) is fast but myopic, committing to locally attractive nodes without ability to backtrack.

SA's advantage stems from its single-solution trajectory search: it evaluates each candidate move in full context (current route, elapsed time, fatigue state, directional costs) without needing to decompose the solution into reusable components.
The Or-opt and swap moves are particularly effective on asymmetric instances because they directly evaluate the actual cost change of rearranging the visit sequence, naturally accounting for directionality.

\paragraph{Parameter tuning.}
The SA parameters were fine-tuned via grid search over a separate set of 1{,000 synthetic instances with randomised control-point distributions across both study areas.}
The tunable parameters and their search ranges were:
\begin{itemize}[nosep]
  \item Initial temperature $T_0 \in \{50, 100, 200, 500\}$
  \item Final temperature $T_{\text{end}} \in \{0.01, 0.1, 1.0\}$
  \item Number of iterations $N_{\text{iter}} \in \{20000, 40000, 80000, 120000\}$
  \item Number of restarts $N_{\text{restart}} \in \{4, 8, 12, 20\}$
  \item Move probabilities $(p_{\text{remove}}, p_{\text{insert}}, p_{\text{or-opt}}, p_{\text{swap}})$, tested in 10 configurations
  \item Or-opt segment length $L_{\text{seg}} \in \{1\text{--}2, 1\text{--}3, 1\text{--}4\}$
  \item Stagnation limit as fraction of $N_{\text{iter}}$: $\{1/3, 1/4, 1/5, 1/6\}$
\end{itemize}

The best configuration, used throughout the experiments, is: $T_0 = 100$, $T_{\text{end}} = 0.1$, move probabilities $(0.30, 0.30, 0.20, 0.20)$, $L_{\text{seg}} \in \{1,2,3\}$, and stagnation limit $N_{\text{iter}}/5$.
The number of iterations and restarts are auto-calibrated based on instance size $n$:
\begin{align}
  N_{\text{iter}} &= \text{clip}(2500n,\; 10000,\; 120000) \\
  N_{\text{restart}} &= \text{clip}(3000/n,\; 40,\; 200)
\end{align}
The rationale is as follows.
The neighbourhood size grows with $n$: more nodes means more possible insertions, removals, and swaps per iteration, so the search space requires proportionally more iterations to explore.
The linear scaling $2500n$ was determined empirically during grid search as the point where additional iterations yielded diminishing returns on solution quality.
The lower bound of $10{,}000$ ensures sufficient exploration even for very small instances, while the upper bound of $120{,}000$ caps runtime.

Conversely, restarts scale inversely with $n$: for small instances ($n \leq 15$), the search space is compact and a single SA run converges quickly, so diversity across restarts (up to 20) is more valuable than depth within a single run.
For large instances ($n \geq 40$), each restart is expensive and the search space is rich enough that a single long run explores effectively, so fewer restarts (minimum 4) suffice.
The total computational effort $N_{\text{iter}} \times N_{\text{restart}}$ remains approximately constant across instance sizes, keeping the SA runtime predictable.

The cooling schedule is geometric: $T_k = T_0 \cdot \rho^k$ with decay rate $\rho = (T_{\text{end}}/T_0)^{1/N_{\text{iter}}}$.

\paragraph{Neighbourhood structure.}
Starting from a greedy insertion solution, the SA applies four neighbourhood moves with the tuned probabilities:
\begin{enumerate}[nosep]
  \item \textbf{Remove} ($p = 0.30$): delete a random node from the route.
  \item \textbf{Insert} ($p = 0.30$): add an unvisited node at its best position (minimum cost increase, computed in $O(k)$ where $k$ is the current route length).
  \item \textbf{Or-opt} ($p = 0.20$): extract a segment of 1--3 nodes and reinsert at a different position.
  \item \textbf{Swap} ($p = 0.20$): exchange a visited node for an unvisited one.
\end{enumerate}

Moves that violate the fatigue budget ($D^f(\sigma') > B$) are rejected.
The search terminates early if no improvement is found for $N_{\text{iter}}/5$ consecutive iterations (stagnation detection).

The full procedure, including the iterated multi-restart wrapper, is given in \cref{alg:sa}.

\begin{algorithm}[H]
\caption{Iterated Simulated Annealing for AOPF}
\label{alg:sa}
\begin{algorithmic}[1]
\Require Cost matrix $\mathbf{C}$, scores $\mathbf{p}$, budget $B$, fatigue rate $\lambda$
\Ensure Best route $\sigma^*$
\State $\sigma^* \gets \emptyset$, $z^* \gets 0$
\State Compute $N_{\text{iter}} \gets \text{clip}(2500n, 10000, 120000)$
\State Compute $N_{\text{restart}} \gets \text{clip}(3000/n, 40, 200)$
\For{$r = 1, \ldots, N_{\text{restart}}$}
  \State $\sigma \gets \textsc{GreedyInsertion}(\mathbf{C}, \mathbf{p}, B)$
  \State $z \gets \sum_{i \in \sigma} p_i$, \; $T \gets T_0$, \; $\rho \gets (T_{\text{end}}/T_0)^{1/N_{\text{iter}}}$, \; $s \gets 0$
  \For{$k = 1, \ldots, N_{\text{iter}}$}
    \If{$s \geq N_{\text{iter}}/5$} \textbf{break} \Comment{Stagnation}
    \EndIf
    \State $T \gets T \cdot \rho$
    \State Draw $u \sim \text{Uniform}(0,1)$
    \If{$u < 0.30$}
      \State $\sigma' \gets \textsc{Remove}(\sigma)$ \Comment{Delete random node}
    \ElsIf{$u < 0.60$}
      \State $\sigma' \gets \textsc{Insert}(\sigma)$ \Comment{Add node at best position}
    \ElsIf{$u < 0.80$}
      \State $\sigma' \gets \textsc{OrOpt}(\sigma)$ \Comment{Relocate segment of 1--3 nodes}
    \Else
      \State $\sigma' \gets \textsc{Swap}(\sigma)$ \Comment{Exchange visited $\leftrightarrow$ unvisited}
    \EndIf
    \If{$D^f(\sigma') > B$}
      \State $s \gets s + 1$; \textbf{continue} \Comment{Infeasible move}
    \EndIf
    \State $z' \gets \sum_{i \in \sigma'} p_i$, \; $\Delta \gets z' - z$
    \If{$\Delta > 0$ \textbf{or} $\exp(\Delta / T) > \text{Uniform}(0,1)$}
      \State $\sigma \gets \sigma'$, \; $z \gets z'$
      \If{$z > z^*$} $z^* \gets z$, \; $\sigma^* \gets \sigma$, \; $s \gets 0$
      \Else \; $s \gets s + 1$
      \EndIf
    \Else \; $s \gets s + 1$
    \EndIf
  \EndFor
\EndFor
\State \Return $\sigma^*$
\end{algorithmic}
\end{algorithm}

\subsection{Branch-and-Cut (B\&C)}
\label{sec:bnc}

The B\&C solver uses the LP relaxation of the AOPF formulation (\cref{sec:formulation}), solved via the HiGHS dual simplex \citep{huangfu2018highs}.
HiGHS was chosen as it is open-source, provides a C API suitable for embedding in a custom B\&C loop, and its dual simplex implementation supports efficient warm-starting after cut addition.
The algorithm proceeds as follows:

\begin{algorithm}[H]
\caption{Branch-and-Cut for AOPF}
\label{alg:bnc}
\begin{algorithmic}[1]
\State Build root LP relaxation
\State Warm-start incumbent $z^* \gets$ SA solution value
\State Push root node onto stack
\While{stack non-empty \textbf{and} time $< T_{\max}$}
  \State Pop node, apply variable fixings
  \For{$k = 1, \ldots, K_{\max}$} \Comment{Cut loop, $K_{\max} = 20$}
    \State Solve LP relaxation
    \If{$\bar{z} \leq z^* + \varepsilon$} \textbf{prune and continue} \EndIf
    \State Find violated SECs via Kosaraju SCC
    \State Find depot-unreachable nodes via reverse BFS
    \State Find violated max-flow connectivity cuts
    \State Find violated budget-based cuts (B1--B4, \cref{sec:separation})
    \If{no violations found} \textbf{break} \EndIf
    \State Add all violated cuts to LP
  \EndFor
  \If{solution is integer-feasible}
    \State Extract route, validate fatigue feasibility
    \If{score $> z^*$} update incumbent \EndIf
  \Else
    \State Branch on most fractional $y_i$
    \State Push child nodes (DFS order)
  \EndIf
\EndWhile
\State \textbf{if} stack empty \textbf{then} solution is provably optimal
\end{algorithmic}
\end{algorithm}

Key implementation details:
\begin{itemize}[nosep]
  \item Branching is performed on the node-visit variables $y_i$ rather than arc variables $x_{ij}$.
        Setting $y_i = 0$ removes node $i$ from all future solutions in the subtree, eliminating $O(n)$ arc variables at once; setting $y_i = 1$ forces the node to be visited.
        By contrast, fixing a single $x_{ij}$ only excludes one arc, requiring up to $O(n^2)$ branching decisions to resolve all fractional arcs.
  \item LP models are deep-copied via \texttt{Highs\_passLp} for thread safety.
  \item Duplicate SECs are tracked via a hash set to avoid redundant cuts.
  \item Maximum branching depth is capped at $D_{\max} = 15$.
  \item Time limit $T_{\max} = 900$\,s (15 minutes).
\end{itemize}

\subsection{Node Selection Strategy}
\label{sec:node-selection}

The B\&C algorithm admits two complementary node selection strategies, and we employ each on the instance class where it is most effective.
\emph{Depth-first search} (DFS) maintains nodes on a stack and always processes the most recently pushed child; it descends quickly into promising subtrees, finds integer-feasible improvements early, and prunes the remaining tree once a strong incumbent is established.
\emph{Best-first search} (BFS) maintains nodes on a priority queue ordered by the LP upper bound and always processes the node most likely to improve the global bound; it minimises the number of nodes examined to close the optimality gap at the cost of slower incumbent discovery.
This tradeoff between depth-first and best-first node selection is a well-studied aspect of branch-and-bound design; \citet{linderoth1999computational} provide a comprehensive computational study of search strategies for mixed integer programming and document that the relative effectiveness of DFS and BFS depends on instance characteristics such as bound quality and tree size.

For the real-terrain instances ($n \leq 40$) we use DFS, and for the benchmark suite (up to $n = 100$) we use BFS.
In preliminary experiments the DFS configuration proved more terrain instances optimal, while the BFS configuration yielded lower mean optimality gaps on the larger benchmark problems.
Both configurations share the same formulation, cut families, and warm-start heuristic.

Both strategies use the same formulation, the same cut families, and the same warm-start heuristic; only the node selection rule and the strong-branching parameters differ.
This allows a clean decomposition of the gains from formulation, cuts, and search strategy, as reported in the results below.

% ════════════════════════════════════════════════════════════════
\section{Computational Pipeline}
\label{sec:pipeline}
% ════════════════════════════════════════════════════════════════

The full pipeline consists of two stages:

\paragraph{Stage 1: Cost Surface Generation (Python).}
\begin{enumerate}[nosep]
  \item Parse OMAP file $\to$ rasterise symbols onto UTM grid $\to$ ACR.
  \item Load DTM, compute TRI, PC, SLO $\to$ normalise $\to$ HCR (\cref{eq:hcr}).
  \item Scale HCR via IQR ratio (\cref{eq:iqr}), merge with ACR (\cref{eq:combined}).
  \item Compute directional slopes from DTM.
  \item Run anisotropic Dijkstra from each node with Minetti factors.
  \item Export asymmetric cost matrix, scores, and budget as JSON.
\end{enumerate}

\paragraph{Stage 2: Route Optimisation (C++).}
\begin{enumerate}[nosep]
  \item Parse JSON input.
  \item Run iterated SA $\to$ warm-start solution.
  \item Build LP model with HiGHS, run B\&C (\cref{alg:bnc}).
  \item Report solution, optimality status, and timing.
\end{enumerate}

% ════════════════════════════════════════════════════════════════
\section{Experimental Setup}
\label{sec:experiments}
% ════════════════════════════════════════════════════════════════

\subsection{Study Areas}

\begin{table}[H]
\centering
\caption{Terrain characteristics of the two study areas.}
\label{tab:terrain}
\begin{tabular}{lcc}
\toprule
& \textbf{Torremocha} & \textbf{La Muela} \\
\midrule
Grid size & $1056 \times 1463$ & $1376 \times 1622$ \\
Resolution & 2.0\,m & 2.0\,m \\
Area & $2.1 \times 2.9$\,km & $2.8 \times 3.2$\,km \\
Elevation range & 447--501\,m & 1164--1628\,m \\
Relief & 55\,m & 464\,m \\
Median slope & $3.2^\circ$ & $15.9^\circ$ \\
Max slope & $58.9^\circ$ & $81.2^\circ$ \\
CRS & EPSG:25829 & EPSG:25830 \\
IQR scaling $c$ & 0.87 & 2.18 \\
HCR median increase & $+14.8\%$ & $+55.1\%$ \\
Cost asymmetry & $16.1\%$ & $50.0\%$ \\
\bottomrule
\end{tabular}
\end{table}

Torremocha is a gently undulating terrain in Extremadura, Spain, with moderate forest cover and extensive rock features.
La Muela is a mountainous area with steep slopes, cliffs, and significant elevation variation.
Additional study areas are being coordinated with the regional orienteering federation to broaden the validation set.

\subsection{Control Distributions}

Seven control-point distributions are generated for each terrain to test solver robustness:
\begin{enumerate}[nosep]
  \item \textbf{Standard}: grid-zone uniform spread (40 controls).
  \item \textbf{Clustered}: 5 hotspot clusters (35 controls).
  \item \textbf{Ring}: concentric rings from depot (30 controls).
  \item \textbf{Path-biased}: 60\% near paths, 40\% far (35 controls).
  \item \textbf{Elevation-biased}: weighted toward high elevation (35 controls).
  \item \textbf{Sparse-far}: few controls, all far from depot (25 controls).
  \item \textbf{Mixed-density}: dense near depot, sparse far (40 controls).
\end{enumerate}

Point values (10--50) are assigned based on terrain difficulty at each control location.

\subsection{Parameters}

Race duration: 1 hour.
Fatigue rate: $\lambda = 0.2$.
Base speed: 2.5\,m/s.
Dijkstra downsampling: $s = 8$.
B\&C time limit: 900\,s.
SA iterations: auto-calibrated per instance.
These values reflect the typical duration of a middle-distance orienteering race.
\paragraph{Calibration under the revised fatigue model (\cref{sec:fatigue}).}
The value $\lambda = 0.2$ above was calibrated for the original model, where $\lambda$ scaled a $[0,1]$-normalised elapsed-budget fraction $t/B$ (grounded in exercise physiology: running economy deteriorates by $\sim$5\% after 60\,min at 80\% $\dot{V}\text{O}_{2\text{max}}$ on flat terrain \citep{sproule1998,zanini2025}, orienteering's 26\% higher oxygen cost than road running \citep{creagh1997} plus $\sim$5\% from uneven surfaces \citep{voloshina2015} approximately doubling that rate, and a further margin for cognitive navigation load \citep{millet2010orienteering}).
Under the revised model (\cref{sec:fatigue}), $\lambda$ multiplies the fatigue state $F_i$, which is expressed in metres of elevation and distance rather than a normalised fraction---the two are not dimensionally comparable, and reusing $\lambda = 0.2$ produces an unphysical result: on a representative synthetic instance, an accumulated state of $F \approx 145\,\mathrm{m}$ gives $1 + \lambda F \approx 30\times$ the base cost, an order of magnitude larger than the $\leq 1.2\times$ ceiling the original model imposed by construction.

We were unable to find a directly usable literature value for $\lambda$ in the units the revised model needs (a pace-decay rate per metre of net elevation gain).
The exercise-physiology and trail-running literature we surveyed reports either whole-race average pace losses (e.g.\ elite ultra-trail runners slowing 15--20\% by race end) or fatigue accumulated as a function of instantaneous work rate/power rather than cumulative elevation---e.g.\ \citet{jaszczak2024} extend Keller's running model to trail racing with a fatigue state that grows proportionally to power output (rate constant $K = 6\times10^{-5}\,\mathrm{s}^{-1}$, taken from prior work rather than fit to their race data)---not a coefficient tying pace decay to cumulative vertical gain specifically.
As a rough order-of-magnitude anchor---not a calibrated or cited constant---we take a representative ultra-trail profile ($\sim$10{,}000\,m cumulative $D^+$, e.g.\ UTMB) together with the $\sim$15--20\% reported end-of-race slowdown: $1 + \lambda D^+ \approx 1.175$ at $D^+ = 10\text{,}000$ gives $\lambda \approx 1.75 \times 10^{-5}$.
Given the absence of a defensible point estimate, we treat $\lambda$ as a swept sensitivity parameter (\texttt{--lambda-sweep} in the solver implementation) over $\{0,\, 1\times10^{-5},\, 1.75\times10^{-5},\, 5\times10^{-5},\, 1\times10^{-4}\}$, centred on the order-of-magnitude anchor above, rather than reporting results at a single unjustified value.

$\rho$ (the downhill-recovery fraction in \cref{eq:psi}) remains a swept sensitivity parameter as before ($\rho \in \{0, 0.25, 0.5, 0.75, 1.0\}$, \texttt{--rho-sweep}): the metabolic-recovery-versus-eccentric-damage trade-off it represents is not settled in the literature we surveyed, so we report sensitivity rather than assert a value.
$\mu$ (the distance-to-elevation weight, \cref{eq:mu}) is the one fatigue parameter derived rather than swept, since it follows directly from the Minetti curve already used for $C_{ij}$; for the terrain parameters in this study it evaluates to $\mu \approx 0.067$ (equivalently, $\sim$14.9 horizontal metres per vertical-metre-equivalent).}

% ════════════════════════════════════════════════════════════════
\section{Results}
\label{sec:results}
% ════════════════════════════════════════════════════════════════

We evaluate the framework on the two real-terrain study areas (Torremocha, 16.1\% cost asymmetry; La Muela, 50.0\% cost asymmetry) and on the synthetic benchmark suite.
For each instance, the SA warm-start solution provides the incumbent; the B\&C then attempts to prove optimality or tighten the LP relaxation gap within a 15-minute time limit.
Both formulations (MTZ and flow) receive the same SA incumbent.
The \emph{gap} reported in each table is the relative difference between the best LP upper bound and the incumbent: $\text{gap} = 100 \times (\bar{z} - z^*) / z^*$, where $\bar{z}$ is the LP bound and $z^*$ is the best integer solution found.
A gap of 0.0\% with ``YES'' indicates that the B\&C has proven the SA solution to be globally optimal.

\subsection{Torremocha Results}

\begin{table}[H]
\centering
\caption{Torremocha results: SA and B\&C under both formulations (IQR scaling $c = 0.87$, cost asymmetry $\mathcal{A} = 16.1\%$). B\&C uses depth-first search (see \cref{sec:node-selection}).}
\label{tab:torremocha}
\begin{tabular}{lrrrrcrrc}
\toprule
& & \multicolumn{2}{c}{\textbf{SA}} & \multicolumn{2}{c}{\textbf{MTZ}} & \multicolumn{2}{c}{\textbf{Flow}} \\
\cmidrule(lr){3-4} \cmidrule(lr){5-6} \cmidrule(lr){7-8}
\textbf{Distribution} & ctrls & pts & visited & gap & optimal? & gap & optimal? \\
\midrule
Standard        & 40 &  940 & 25 & 14.1\% & no  & 12.2\% & no \\
Clustered       & 35 &  510 & 29 &  0.0\% & \textbf{YES} &  0.0\% & \textbf{YES} \\
Ring            & 30 &  650 & 26 &  0.0\% & \textbf{YES} &  0.0\% & \textbf{YES} \\
Path-biased     & 35 &  630 & 30 &  0.0\% & \textbf{YES} &  0.0\% & \textbf{YES} \\
Elev-biased     & 35 &  380 & 29 &  7.4\% & no  &  6.3\% & no \\
Sparse-far      & 25 &  660 & 19 &  0.0\% & \textbf{YES} &  0.0\% & \textbf{YES} \\
Mixed-density   & 40 &  680 & 32 &  8.6\% & no  &  7.1\% & no \\
\bottomrule
\end{tabular}
\end{table}

The flow formulation proves 4 of 7 instances optimal, matching MTZ (4/7) on the same subset (Clustered, Ring, Path-biased, Sparse-far).
The flow formulation produces tighter gaps on the unproven instances, with improvements on Standard (14.1\% $\to$ 12.2\%), Elev-biased (7.4\% $\to$ 6.3\%), and Mixed-density (8.6\% $\to$ 7.1\%).

\subsection{La Muela Results}

\begin{table}[H]
\centering
\caption{La Muela results: SA and B\&C under both formulations (IQR scaling $c = 2.18$, cost asymmetry $\mathcal{A} = 50.0\%$). B\&C uses depth-first search (see \cref{sec:node-selection}).}
\label{tab:lamuela}
\begin{tabular}{lrrrrcrrc}
\toprule
& & \multicolumn{2}{c}{\textbf{SA}} & \multicolumn{2}{c}{\textbf{MTZ}} & \multicolumn{2}{c}{\textbf{Flow}} \\
\cmidrule(lr){3-4} \cmidrule(lr){5-6} \cmidrule(lr){7-8}
\textbf{Distribution} & ctrls & pts & visited & gap & optimal? & gap & optimal? \\
\midrule
Standard        & 31 &  610 & 14 &  0.0\% & \textbf{YES} &  0.0\% & \textbf{YES} \\
Clustered       & 35 &  800 & 24 &  0.0\% & \textbf{YES} &  0.0\% & \textbf{YES} \\
Ring            & 30 &  470 & 21 &  6.3\% & no  &  0.0\% & \textbf{YES} \\
Path-biased     & 35 &  570 & 17 &  0.0\% & \textbf{YES} &  0.0\% & \textbf{YES} \\
Elev-biased     & 35 &  730 & 22 & 17.0\% & no  & 12.4\% & no \\
Sparse-far      & 25 &  600 & 13 &  0.0\% & \textbf{YES} &  0.0\% & \textbf{YES} \\
Mixed-density   & 40 &  680 & 25 & 14.6\% & no  &  8.6\% & no \\
\bottomrule
\end{tabular}
\end{table}

On La Muela, the flow formulation proves 5 of 7 instances optimal versus 4 for MTZ, with the flow formulation additionally proving Ring optimal.
The flow formulation reduces gaps on the unproven instances: Elev-biased (17.0\% $\to$ 12.4\%) and Mixed-density (14.6\% $\to$ 8.6\%).
Despite the significantly harder terrain (464\,m relief, 50\% cost asymmetry), the flow-based B\&C closes the optimality gap substantially on all instances.

Across both terrains (14 instances total), the flow formulation proves 9 of 14 instances optimal (vs.\ 8 for MTZ) and reduces the mean gap on unproven instances from 11.3\% to 9.3\%.

\subsection{Synthetic Benchmark Instances}
\label{sec:benchmark}

To the best of our knowledge, no standard benchmark exists for the asymmetric OP with fatigue.
We generate a parameterized benchmark suite by varying four dimensions independently from a baseline configuration ($n = 40$, $\alpha = 0.25$, $\lambda = 0.2$, medium budget):
\begin{itemize}[nosep]
  \item Node count $n \in \{20, 30, 40, 50, 75, 100\}$
  \item Asymmetry parameter $\alpha \in \{0.0, 0.1, 0.25, 0.5\}$, controlling directional cost perturbation
  \item Fatigue rate $\lambda \in \{0.0, 0.1, 0.2, 0.3\}$
  \item Budget tightness: loose ($\sim$70\%), medium ($\sim$50\%), tight ($\sim$30\%), defined as a fraction of the nearest-neighbour tour cost
\end{itemize}
Additionally, four extreme-case instances test boundary conditions (e.g., symmetric with no fatigue, highly asymmetric with tight budget).
Nodes are placed uniformly at random in a $1000 \times 1000$\,m domain with the depot at the centre.
Asymmetric costs are generated by perturbing Euclidean distances with a directional slope factor proportional to $\alpha$, plus 10\% random terrain noise.
Scores are correlated with distance from the depot (10--50 points, rounded to multiples of 10).
All instances and the solver are released as part of the benchmark suite.
Note that the synthetic asymmetry is spatially uncorrelated, unlike real terrain where slope-induced cost differences exhibit spatial autocorrelation; the real-terrain experiments on Torremocha and La Muela complement the benchmark by testing the solver on naturally structured asymmetry.

\subsection{Benchmark Results}

Both formulations (MTZ and flow) are evaluated on the full benchmark suite under identical conditions (same SA warm start, 900\,s time limit).

\begin{table}[H]
\centering
\caption{Benchmark results: effect of instance size ($\alpha = 0.25$, $\lambda = 0.2$, medium budget). Both formulations receive the same SA warm start.}
\label{tab:bench_size}
\begin{tabular}{rrrrcrrc}
\toprule
& & & \multicolumn{2}{c}{\textbf{MTZ}} & \multicolumn{2}{c}{\textbf{Flow}} \\
\cmidrule(lr){4-5} \cmidrule(lr){6-7}
$n$ & SA pts & visited & gap & optimal? & gap & optimal? \\
\midrule
20  &  350 & 12 &  0.0\% & \textbf{YES} &  0.0\% & \textbf{YES} \\
30  &  540 & 18 &  0.0\% & \textbf{YES} &  0.0\% & \textbf{YES} \\
40  &  660 & 26 & 18.5\% & no &  7.5\% & no \\
50  &  940 & 31 & 16.5\% & no &  5.9\% & no \\
75  & 1370 & 43 & 19.4\% & no &  6.5\% & no \\
100 & 1840 & 63 & 19.9\% & no &  7.8\% & no \\
\bottomrule
\end{tabular}
\end{table}

\begin{table}[H]
\centering
\caption{Benchmark results: effect of asymmetry and fatigue ($n = 40$, medium budget).}
\label{tab:bench_params}
\begin{tabular}{llrrcrrc}
\toprule
& & & \multicolumn{2}{c}{\textbf{MTZ}} & \multicolumn{2}{c}{\textbf{Flow}} \\
\cmidrule(lr){4-5} \cmidrule(lr){6-7}
Parameter & Value & SA pts & gap & optimal? & gap & optimal? \\
\midrule
\multirow{4}{*}{Asymmetry $\alpha$}
& 0.00 &  690 & 15.8\% & no &  6.2\% & no \\
& 0.10 &  730 & 14.1\% & no &  2.8\% & no \\
& 0.25 &  670 & 20.2\% & no &  5.9\% & no \\
& 0.50 &  630 & 22.7\% & no &  7.5\% & no \\
\midrule
\multirow{4}{*}{Fatigue $\lambda$}
& 0.00 &  790 &  6.3\% & no &  0.0\% & \textbf{YES} \\
& 0.10 &  730 &  8.5\% & no &  3.0\% & no \\
& 0.20 &  730 & 15.4\% & no &  7.7\% & no \\
& 0.30 &  750 & 30.5\% & no & 10.5\% & no \\
\bottomrule
\end{tabular}
\end{table}

The benchmark reveals several patterns.
First, the flow formulation consistently produces tighter LP bounds than MTZ across instance sizes, with the gap reduction growing from 11.0\,pp at $n = 40$ to 12.1\,pp at $n = 100$ (\cref{tab:bench_size}).
Second, the flow formulation reduces asymmetry-related gaps substantially: from 22.7\% (MTZ) to 7.5\% (flow) at $\alpha = 0.5$ (\cref{tab:bench_params}).
Third, the fatigue dimension confirms the McCormick looseness under MTZ: the MTZ gap grows from 6.3\% ($\lambda = 0$) to 30.5\% ($\lambda = 0.3$), while the flow formulation shows a much smaller increase from 0.0\% to 10.5\%.
Across all 21 instances, the flow formulation reduces the mean gap from 16.0\% to 5.5\% and proves 7 instances optimal versus 4 for MTZ, with a maximum gap of 10.5\% compared to 42.0\% for MTZ.

\begin{table}[H]
\centering
\caption{Benchmark results: budget tightness and extreme cases ($n = 40$ unless noted).}
\label{tab:bench_extra}
\begin{tabular}{lrrcrrc}
\toprule
& & \multicolumn{2}{c}{\textbf{MTZ}} & \multicolumn{2}{c}{\textbf{Flow}} \\
\cmidrule(lr){3-4} \cmidrule(lr){5-6}
Instance & SA pts & gap & optimal? & gap & optimal? \\
\midrule
Loose budget          & 960  & 18.2\% & no  &  8.0\% & no \\
Medium budget         & 680  & 25.5\% & no  &  5.3\% & no \\
Tight budget          & 450  &  0.0\% & \textbf{YES} &  0.0\% & \textbf{YES} \\
Symmetric easy ($n{=}20$, $\alpha{=}0$, $\lambda{=}0$) & 470 & 0.0\% & \textbf{YES} & 0.0\% & \textbf{YES} \\
Asymmetric hard ($n{=}100$, $\alpha{=}0.5$, $\lambda{=}0.3$) & 1160 & 42.0\% & no & 10.5\% & no \\
High asym., no fatigue ($\alpha{=}0.5$, $\lambda{=}0$) & 970 & 9.1\% & no & 1.1\% & no \\
Symmetric, high fatigue ($\alpha{=}0$, $\lambda{=}0.3$) & 850 & 24.9\% & no & 9.9\% & no \\
\bottomrule
\end{tabular}
\end{table}

The extreme cases highlight the flow formulation's strengths: the hardest instance (asymmetric hard) improves from 42.0\% to 10.5\%, and the high-asymmetry no-fatigue instance drops from 9.1\% to 1.1\%.
The medium-budget instance shows a large improvement (25.5\% $\to$ 5.3\%), while both formulations prove the tight-budget and symmetric-easy instances optimal.
The no-fatigue instance ($\lambda = 0$) is newly proven optimal by the flow formulation.

\begin{proposition}[Optimality certificate]
\label{prop:optimality}
Let $z^*_{\text{SA}}$ denote the objective value of the SA incumbent solution, and let $\mathcal{T}$ denote the Branch-and-Cut search tree.
For each node $\nu \in \mathcal{T}$, let $\bar{z}(\nu)$ denote the optimal value of the LP relaxation at $\nu$ (an upper bound on the best integer-feasible solution in the subtree rooted at $\nu$).
If the B\&C algorithm terminates with an empty node stack (i.e., all nodes have been either explored or pruned), then $z^*_{\text{SA}}$ is the globally optimal value of the AOPF.
\end{proposition}

\begin{proof}
The B\&C algorithm partitions the feasible region of the AOPF into a binary search tree $\mathcal{T}$ by branching on fractional $y_i$ variables.
At each node $\nu$, the LP relaxation (with dynamically added subtour elimination and connectivity cuts) yields an upper bound $\bar{z}(\nu) \geq z^*(\nu)$, where $z^*(\nu)$ is the optimal integer solution in the subtree rooted at $\nu$.
This bound is valid because:
\begin{enumerate}[nosep]
  \item The LP relaxation is a relaxation of the integer program (all integer constraints are dropped, and the feasible set is enlarged).
  \item The subtour elimination cuts (\cref{eq:sec}) and connectivity cuts (\cref{eq:outcut,eq:incut}) are valid inequalities---they do not exclude any integer-feasible solution.
\end{enumerate}

A node $\nu$ is pruned if $\bar{z}(\nu) \leq z^*_{\text{SA}} + \varepsilon$ (with $\varepsilon = 10^{-6}$), meaning no integer solution in its subtree can improve upon the incumbent.
A node is also pruned if its LP relaxation is infeasible, or if an integer-feasible solution is found and used to update the incumbent.

When the algorithm terminates with an empty stack, every leaf of $\mathcal{T}$ has been resolved by one of:
\begin{itemize}[nosep]
  \item \emph{Bound pruning}: $\bar{z}(\nu) \leq z^*_{\text{SA}} + \varepsilon$, certifying that no better solution exists in this subtree.
  \item \emph{Infeasibility}: the LP relaxation at $\nu$ has no feasible solution.
  \item \emph{Integer solution}: an integer-feasible solution was found and, if superior, became the new incumbent.
\end{itemize}

Since the branching is exhaustive (every fractional variable is eventually fixed to 0 or 1) and the union of all subtrees covers the entire feasible region of the AOPF, the incumbent at termination is a globally optimal solution.
In particular, if the incumbent remains the SA solution $z^*_{\text{SA}}$, then no integer-feasible solution with objective value strictly greater than $z^*_{\text{SA}}$ exists.
\end{proof}

\subsection{Ablation Study: Effect of Cut Families}
\label{sec:ablation}

To quantify the marginal contribution of each cut family introduced in \cref{sec:separation}, we evaluate five cumulative configurations of the flow-based B\&C solver on the full benchmark suite of 21 instances, following the component evaluation methodology of \citet{kobeaga2021revisited}.
Each configuration adds one cut family to the previous one, using the same SA warm start and 900\,s time limit.
\Cref{tab:ablation} reports the mean optimality gap, number of instances proven optimal, and mean solving time for each configuration.

\begin{table}[H]
\centering
\caption{Ablation study on the benchmark suite (21 instances, 900\,s time limit). Each row adds one cut family to the previous configuration. B\&C uses best-first search (see \cref{sec:node-selection}). Gap is $100 \times (\bar{z} - z^*) / z^*$.}
\label{tab:ablation}
\small
\begin{tabular}{clrrrr}
\toprule
 & \textbf{Cuts enabled} & \textbf{Mean} & \textbf{Max} & \textbf{Opt.} & \textbf{Mean} \\
 &  & \textbf{gap} & \textbf{gap} &  & \textbf{time (s)} \\
\midrule
1 & Base (SECs + conn.\ + tight.\ + fat.)  & 5.54\% & 20.32\% & 6/21 & 727 \\
2 & + B1 (lifted covers)                    & 5.85\% & 19.88\% & 7/21 & 741 \\
3 & + B2 (routing infeas.)                  & 5.41\% & 19.69\% & 7/21 & 723 \\
4 & + B3 (cycle covers)                     & 5.35\% & 20.19\% & 7/21 & 718 \\
5 & + B4 (path ineq.)                       & 5.42\% & 19.81\% & 7/21 & 718 \\
\bottomrule
\end{tabular}
\end{table}

The base configuration---comprising subtour elimination, connectivity cuts, tightened flow-arc coupling, and fatigue-aware cover weighting---proves 6 of 21 instances optimal with a mean gap of 5.54\%.
Adding B1 (lifted cover cuts) proves one additional instance (bench\_020, high asymmetry, no fatigue) but slightly increases the mean gap to 5.85\% due to the overhead of cover separation on harder instances.
B2 (routing infeasibility cuts) reduces the mean gap to 5.41\% without changing the number of optimal solutions, confirming that the Hungarian-based separation tightens bounds on non-optimal instances.
B3 (cycle cover cuts) yields the lowest mean gap (5.35\%) and the fastest mean time on bench\_020 (690\,s vs.\ 838\,s with B1+B2 alone).
B4 (path inequalities) provides no further improvement: the mean gap increases slightly to 5.42\% and bench\_020 narrowly misses the time limit (gap $< 0.1\%$), suggesting that the additional LP overhead outweighs the tightening effect.

% ════════════════════════════════════════════════════════════════
\section{Discussion}
\label{sec:discussion}
% ════════════════════════════════════════════════════════════════

\paragraph{SA effectiveness.}
On all 9 instances where the B\&C proves optimality, the incumbent at termination is the original SA warm-start solution---the B\&C confirms the SA solution is optimal but does not improve upon it.
This suggests SA is highly effective for OP instances of this size ($n \leq 40$), consistent with the literature on metaheuristics for the OP \citep{vansteenwegen2011orienteering}.

\paragraph{Cost asymmetry.}
The 50\% asymmetry on La Muela demonstrates that symmetric OP formulations would be inadequate for mountainous terrain.
The Minetti model captures this directionality naturally through the metabolic cost polynomial.

\paragraph{Fatigue linearisation.}
The comparison between the MTZ and flow formulations confirms that the McCormick envelope is the primary source of LP relaxation looseness in the MTZ approach.
The flow formulation, combined with max-flow separation and reliability branching, reduces the mean benchmark gap from 13.2\% to 6.3\% and halves the maximum gap from 22.7\% to 15.2\%.
On real-terrain instances, the flow formulation reduces the La Muela Mixed-density gap from 14.6\% (MTZ) to 8.6\%, demonstrating the practical impact of this tighter relaxation.

\paragraph{Computational complexity and runtime.}
The overall framework has three computational bottlenecks, each with distinct scaling behaviour.

\emph{Cost matrix construction.}
For $n$ control points on a grid of $H \times W$ cells downsampled by factor $s$, the anisotropic Dijkstra runs $n$ times on a grid of $(H/s)(W/s)$ cells.
Each Dijkstra has complexity $O\bigl(\frac{HW}{s^2} \log \frac{HW}{s^2}\bigr)$, giving a total cost matrix construction time of $O\bigl(\frac{nHW}{s^2} \log \frac{HW}{s^2}\bigr)$.
In practice, with Numba JIT compilation and $s = 8$, the cost matrix for 41 nodes on a $1056 \times 1463$ grid computes in under 4 seconds.

\emph{Simulated Annealing.}
Each SA iteration evaluates a neighbourhood move in $O(k)$ time, where $k$ is the current route length.
With $N_{\text{iter}}$ iterations and $N_{\text{restart}}$ restarts, the total SA time is $O(N_{\text{restart}} \cdot N_{\text{iter}} \cdot k)$.
Since $k \leq n$ and both $N_{\text{iter}}$ and $N_{\text{restart}}$ are bounded by constants dependent on $n$, the SA runs in $O(n^2)$ per restart in the worst case.
Empirically, SA completes in under 1 second for all tested instances ($n \leq 41$).

\emph{Branch-and-Cut.}
The B\&C has worst-case exponential complexity in $n$, as the number of branching nodes can grow as $O(2^n)$.
The LP relaxation at each node has $O(n^2)$ variables and $O(n^2)$ constraints (before cuts), solvable in polynomial time by the simplex method.
However, the number of subtour elimination cuts can be exponential in $|V|$.
In our experiments, the B\&C runtime varies dramatically with instance structure:
instances with fewer nodes or simpler topology (e.g., Clustered with 24 nodes: 14.1\,s) are solved orders of magnitude faster than dense instances (e.g., Standard with 31 nodes: 900\,s, hitting the time limit).
This exponential sensitivity to instance structure is characteristic of branch-and-cut methods for routing problems and motivates the SA warm start, which provides a strong incumbent that enables aggressive pruning early in the search.

\paragraph{Limitations.}
The framework has not been validated against GPS tracks from real orienteers.
The fatigue model is linear, whereas real fatigue may follow a more complex curve (e.g., exponential degradation or threshold effects at high intensity).
The B\&C does not scale well beyond $\sim$50 nodes due to the exponential growth of the branching tree; for larger instances, the SA solution should be treated as a high-quality heuristic without an optimality guarantee.
The Dijkstra downsampling ($s = 8$) introduces path approximation error that has not been formally quantified, though empirical comparison with $s = 4$ suggests the effect on the cost matrix is small ($< 2\%$ median deviation).

% ════════════════════════════════════════════════════════════════
\section{Conclusion}
\label{sec:conclusion}
% ════════════════════════════════════════════════════════════════
 
We presented an integrated framework for optimal route planning in orienteering, combining GIS-based cost surface generation with exact combinatorial optimisation.
The key methodological contributions are:
(i)~a GIS pipeline that transforms orienteering map data and terrain models into a fully asymmetric, fatigue-aware cost matrix, naturally giving rise to the Asymmetric Orienteering Problem with Fatigue (AOPF);
(ii)~IQR-based raster fusion adapting to terrain characteristics;
(iii)~two Branch-and-Cut formulations for the AOPF---an MTZ-based model with McCormick linearisation and a flow-based model with a natively linear fatigue budget;
(iv)~valid inequalities tailored to the AOPF: tightened flow-arc coupling, fatigue-aware lifted covers, and routing infeasibility cuts (a novel class based on the assignment relaxation), complemented by directed adaptations of cycle cover and path inequalities;
(v)~a parameterised benchmark suite for the asymmetric OP with fatigue; and
(vi)~computational evidence that Simulated Annealing finds optimal solutions on real-terrain instances, with an ablation study quantifying the contribution of each cut family.

Future work includes validation against real GPS data, extension to multi-day events, and integration of weather and time-of-day effects on terrain traversability.

\paragraph{Reproducibility.}
The source code, benchmark instances (14 JSON files with asymmetric cost matrices and proven optimal solutions), and supplementary data are publicly available at \url{https://github.com/PabloAlmendralejo/Orienteering-Problem/tree/main}.
To the best of our knowledge, no benchmark instances exist for the asymmetric OP with terrain-based costs and fatigue; we release ours to facilitate future comparison.

% ════════════════════════════════════════════════════════════════
% REFERENCES
% ════════════════════════════════════════════════════════════════

\bibliographystyle{plainnat}

\begin{thebibliography}{20}

\bibitem[Tsiligirides(1984)]{tsiligirides1984heuristic}
T.~Tsiligirides.
\newblock Heuristic methods applied to orienteering.
\newblock \emph{Journal of the Operational Research Society}, 35(9):797--809, 1984.

\bibitem[Minetti et~al.(2002)]{minetti2002}
A.~E. Minetti, C.~Moia, G.~S. Roi, D.~Susta, and G.~Ferretti.
\newblock Energy cost of walking and running at extreme uphill and downhill slopes.
\newblock \emph{Journal of Applied Physiology}, 93(3):1039--1046, 2002.

\bibitem[Albert and S\'{a}rk\"{o}zy(2021)]{ica2021lcp}
G.~Albert and Z.~S\'{a}rk\"{o}zy.
\newblock Route planning on orienteering maps with least-cost path analysis.
\newblock \emph{Proceedings of the ICA}, 4:4, 2021.
\newblock \doi{10.5194/ica-proc-4-4-2021}.

\bibitem[Riley et~al.(1999)]{riley1999tri}
S.~J. Riley, S.~D. DeGloria, and R.~Elliot.
\newblock A terrain ruggedness index that quantifies topographic heterogeneity.
\newblock \emph{Intermountain Journal of Sciences}, 5(1--4):23--27, 1999.

\bibitem[Zevenbergen and Thorne(1987)]{zevenbergen1987}
L.~W. Zevenbergen and C.~R. Thorne.
\newblock Quantitative analysis of land surface topography.
\newblock \emph{Earth Surface Processes and Landforms}, 12(1):47--56, 1987.

\bibitem[Huangfu and Hall(2018)]{huangfu2018highs}
Q.~Huangfu and J.~A.~J. Hall.
\newblock Parallelizing the dual revised simplex method.
\newblock \emph{Mathematical Programming Computation}, 10(1):119--142, 2018.

\bibitem[Vansteenwegen et~al.(2011)]{vansteenwegen2011orienteering}
P.~Vansteenwegen, W.~Souffriau, and D.~Van~Oudheusden.
\newblock The orienteering problem: A survey.
\newblock \emph{European Journal of Operational Research}, 209(1):1--10, 2011.

\bibitem[Fischetti et~al.(1998)]{fischetti1998solving}
M.~Fischetti, J.~J. Salazar~Gonz\'{a}lez, and P.~Toth.
\newblock Solving the orienteering problem through branch-and-cut.
\newblock \emph{INFORMS Journal on Computing}, 10(2):133--148, 1998.

\bibitem[Gendreau et~al.(1998)]{gendreau1998generalized}
M.~Gendreau, G.~Laporte, and F.~Semet.
\newblock A branch-and-cut algorithm for the undirected selective travelling salesman problem.
\newblock \emph{Networks}, 32(4):263--273, 1998.

\bibitem[Gunawan et~al.(2016)]{gunawan2016orienteering}
A.~Gunawan, H.~C. Lau, and P.~Vansteenwegen.
\newblock Orienteering problem: A survey of recent variants, solution approaches and applications.
\newblock \emph{European Journal of Operational Research}, 255(2):315--332, 2016.

\bibitem[Kirkpatrick et~al.(1983)]{kirkpatrick1983sa}
S.~Kirkpatrick, C.~D. Gelatt, and M.~P. Vecchi.
\newblock Optimization by simulated annealing.
\newblock \emph{Science}, 220(4598):671--680, 1983.

\bibitem[Tasgetiren(2002)]{tasgetiren2002ga_op}
M.~F. Tasgetiren.
\newblock A genetic algorithm with an adaptive penalty function for the orienteering problem.
\newblock \emph{Journal of the Operational Research Society}, 53(3):263--270, 2002.

\bibitem[Ke et~al.(2008)]{ke2008aco_op}
L.~Ke, C.~Archetti, and Z.~Feng.
\newblock Ants can solve the team orienteering problem.
\newblock \emph{Computers \& Industrial Engineering}, 54(3):648--665, 2008.

\bibitem[Vansteenwegen et~al.(2009)]{vansteenwegen2009iterated}
P.~Vansteenwegen, W.~Souffriau, G.~Vanden~Berghe, and D.~Van~Oudheusden.
\newblock Iterated local search for the team orienteering problem with time windows.
\newblock \emph{Computers \& Operations Research}, 36(12):3281--3290, 2009.

\bibitem[Tobler(1993)]{tobler1993}
W.~Tobler.
\newblock Three presentations on geographical analysis and modeling.
\newblock Technical Report 93-1, National Center for Geographic Information and Analysis, 1993.

\bibitem[Rees(2004)]{rees2004}
W.~G. Rees.
\newblock Least-cost paths in mountainous terrain.
\newblock \emph{Computers \& Geosciences}, 30(3):203--209, 2004.

\bibitem[Sproule(1998)]{sproule1998}
J.~Sproule.
\newblock Running economy deteriorates following 60 min of exercise at 80\% {VO2max}.
\newblock \emph{European Journal of Applied Physiology}, 77:366--371, 1998.

\bibitem[Zanini et~al.(2025)]{zanini2025}
M.~Zanini, R.~C. Blagrove, and J.~P. Folland.
\newblock The effect of 90 and 120 min of running on the determinants of endurance performance in well-trained male marathon runners.
\newblock \emph{Scandinavian Journal of Medicine \& Science in Sports}, 35(6):e70042, 2025.

\bibitem[Creagh and Reilly(1997)]{creagh1997}
U.~Creagh and T.~Reilly.
\newblock Physiological and biomechanical aspects of orienteering.
\newblock \emph{Sports Medicine}, 24(6):409--418, 1997.

\bibitem[Voloshina et~al.(2015)]{voloshina2015}
A.~S. Voloshina, A.~D. Kuo, M.~A. Daley, and D.~P. Ferris.
\newblock Biomechanics and energetics of running on uneven terrain.
\newblock \emph{Journal of Experimental Biology}, 218(5):711--719, 2015.

\bibitem[Millet et~al.(2010)]{millet2010orienteering}
G.~Y. Millet, C.~Divert, T.~Banizette, and J.-B. Morin.
\newblock Changes in running pattern due to fatigue and cognitive load in orienteering.
\newblock \emph{Journal of Sports Sciences}, 28(2):153--160, 2010.

\bibitem[Jaszczak and P{\l}ociniczak(2024)]{jaszczak2024}
B.~Jaszczak and {\L}.~P{\l}ociniczak.
\newblock Optimal strategy for trail running with nutrition and fatigue factors.
\newblock \emph{arXiv preprint arXiv:2401.02919}, 2024.

\bibitem[Fischetti(1991)]{fischetti1991atsp}
M.~Fischetti.
\newblock Facets of the asymmetric traveling salesman polytope.
\newblock \emph{Mathematics of Operations Research}, 16(1):42--56, 1991.

\bibitem[Fischetti and Toth(1997)]{fischetti1997atsp}
M.~Fischetti and P.~Toth.
\newblock A polyhedral approach to the asymmetric traveling salesman problem.
\newblock \emph{Management Science}, 43(11):1520--1536, 1997.

\bibitem[Bekta\c{s} and Laporte(2011)]{bektas2011pollution}
T.~Bekta\c{s} and G.~Laporte.
\newblock The pollution-routing problem.
\newblock \emph{Transportation Research Part B: Methodological}, 45(8):1232--1250, 2011.

\bibitem[Malandraki and Daskin(1992)]{malandraki1992timedep}
C.~Malandraki and M.~S. Daskin.
\newblock Time dependent vehicle routing problems: Formulations, properties and heuristic algorithms.
\newblock \emph{Transportation Science}, 26(3):185--200, 1992.

\bibitem[Kara et~al.(2007)]{kara2007energy}
I.~Kara, B.~Y. Kara, and M.~K. Yetis.
\newblock Energy minimizing vehicle routing problem.
\newblock In \emph{Combinatorial Optimization and Applications (COCOA)}, pages 62--71. Springer, 2007.

\bibitem[Feillet et~al.(2005)]{feillet2005tsp_profits}
D.~Feillet, P.~Dejax, and M.~Gendreau.
\newblock Traveling salesman problems with profits.
\newblock \emph{Transportation Science}, 39(2):188--205, 2005.

\bibitem[Kobeaga et~al.(2021)]{kobeaga2021revisited}
G.~Kobeaga, M.~Merino, and J.~A. Lozano.
\newblock A revisited branch-and-cut algorithm for large-scale orienteering problems.
\newblock \emph{European Journal of Operational Research}, 290(3):865--880, 2021.

\bibitem[Linderoth and Savelsbergh(1999)]{linderoth1999computational}
J.~T. Linderoth and M.~W.~P. Savelsbergh.
\newblock A computational study of search strategies for mixed integer programming.
\newblock \emph{INFORMS Journal on Computing}, 11(2):173--187, 1999.

\end{thebibliography}

\end{document}
